%  LaTeX support: latex@mdpi.com 
%  For support, please attach all files needed for compiling as well as the log file, and specify your operating system, LaTeX version, and LaTeX editor.

%=================================================================
\documentclass[sustainability,article,accept,pdftex,moreauthors]{Definitions/mdpi} 

%=================================================================
% MDPI internal commands
\firstpage{1} 
\makeatletter 
\setcounter{page}{\@firstpage} 
\makeatother
\pubvolume{1}
\issuenum{1}
\articlenumber{0}
\pubyear{2022}
\copyrightyear{2022}
\externaleditor{Academic Editors: Chong Xu and Jian Chen
}
\datereceived{4 November 2022} 
\dateaccepted{26 November 2022} 
\datepublished{} 
%\datecorrected{} % Corrected papers include a "Corrected: XXX" date in the original paper.
%\dateretracted{} % Corrected papers include a "Retracted: XXX" date in the original paper.
\hreflink{https://doi.org/} % If needed use \linebreak
\graphicspath{{figures/}} % path to folder with figures
%\doinum{}

\usepackage{listings}
\usepackage{xcolor}

\definecolor{codegreen}{rgb}{0,0.6,0}
\definecolor{codegray}{rgb}{0.5,0.5,0.5}
\definecolor{codepurple}{RGB}{204, 0, 153}
\definecolor{backcolour}{RGB}{255, 250, 230}
%    
\lstdefinestyle{mystyle}{
    backgroundcolor=\color{backcolour}, 
    commentstyle=\color{codegreen},
    keywordstyle=\color{magenta},
    numberstyle=\tiny\color{codegray},
    stringstyle=\color{codepurple},
    basicstyle=\ttfamily\scriptsize,
    breakatwhitespace=false,         
    breaklines=true,                 
    captionpos=b,                    
    keepspaces=true,                 
    numbers=left,                    
    numbersep=5pt,                  
    showspaces=false,                
    showstringspaces=false,
    showtabs=false,                  
    tabsize=2
}

\lstset{style=mystyle}
	
\usepackage{soul} % highlighting text
%------------------------------------------------------------------
% The following line should be uncommented if the LaTeX file is uploaded to arXiv.org
%\pdfoutput=1

%=================================================================
% Add packages and commands here. The following packages are loaded in our class file: fontenc, inputenc, calc, indentfirst, fancyhdr, graphicx, epstopdf, lastpage, ifthen, lineno, float, amsmath, setspace, enumitem, mathpazo, booktabs, titlesec, etoolbox, tabto, xcolor, soul, multirow, microtype, tikz, totcount, changepage, attrib, upgreek, cleveref, amsthm, hyphenat, natbib, hyperref, footmisc, url, geometry, newfloat, caption

%=================================================================
%% Please use the following mathematics environments: Theorem, Lemma, Corollary, Proposition, Characterization, Property, Problem, Example, ExamplesandDefinitions, Hypothesis, Remark, Definition, Notation, Assumption
%% For proofs, please use the proof environment (the amsthm package is loaded by the MDPI class).

%=================================================================
% Full title of the paper (Capitalized)

\Title{Seismotectonics %MDPI: 
 of Shallow-Focus Earthquakes in Venezuela with Links to Gravity Anomalies and Geologic Heterogeneity Mapped by a GMT Scripting Language} % Please confirm this change to the title <-- Answer: yes, we confirm.

% MDPI internal command: Title for citation in the left column
\TitleCitation{Seismotectonics of Shallow-Focus Earthquakes in Venezuela with Links to Gravity Anomalies and Geologic Heterogeneity Mapped by a GMT Scripting Language}

% Author Orchid ID: enter ID or remove command
\newcommand{\orcidauthorA}{0000-0002-5759-1089} % Add \orcidA{} behind the author's name
\newcommand{\orcidauthorB}{0000-0002-6461-1551} % Add \orcidB{} behind the author's name

% Authors, for the paper (add full first names)
\Author{Polina Lemenkova %MDPI: Please carefully check the accuracy of names and affiliations. <-- Answer: yes, we confirm that names and affiliations are correct.
 *\orcidA{} and Olivier Debeir \orcidB{}}

% MDPI internal command: Authors, for metadata in PDF
\AuthorNames{Poina Lemenkova and Olivier Debeir}

% MDPI internal command: Authors, for citation in the left column
\AuthorCitation{Lemenkova, P.; Debeir, O.}
% If this is a Chicago style journal: Lastname, Firstname, Firstname Lastname, and Firstname Lastname.

% Affiliations / Addresses (Add [1] after \address if there is only one affiliation.)
\address[1]{ Laboratory of Image Synthesis and Analysis (LISA), École Polytechnique de Bruxelles (Brussels Faculty of Engineering), %MDPI: please confirm if it is correct to put it in brackets. <-- Answer: yes, correct (first official name, then translation into English in brackets).
 Université Libre de Bruxelles (ULB), %MDPI: The provided information should be arranged from subordinate to superior, please confirm <-- Answer: yes, correct (first Laboratory, then Faculty, finally University)
Building L, Campus de Solbosch, ULB---LISA CP165/57, Avenue Franklin D. Roosevelt 50, 1000 Brussels, Belgium; olivier.debeir@ulb.be %MDPI: We added the email addresses here according to the submitting system. Please confirm.   %MDPI: This is an important note to let you know that to protect the privacy of the author's contact information, we will only display the corresponding authors' contact information on the published paper. However, we hope the contact information of the other authors can also be confirmed or corrected during the proofreading stage, which will be recorded in our database to furnish the author's integrated publication history and also contribute to our future communications. <-- Answer: yes, the emails are correct. We agree with this information.

}

% Contact information of the corresponding author
\corres{\hangafter=1 \hangindent=1.05em \hspace{-0.82em}Correspondence:  polina.lemenkova@ulb.be; Tel.: +32-471860459}

% The commands \thirdnote{} till \eighthnote{} are available for further notes

%\simplesumm{} % Simple summary

% Abstract (Do not insert blank lines, i.e. \\) 
\abstract{This paper presents a cartographic framework based on algorithms of GMT codes for mapping seismically active areas in Venezuela. The data included raster grids from GEBCO, EGM-2008, and vector geological layers from the USGS. The data were  iteratively processed in the console of GMT, converted by GDAL, formatted, and mapped for geophysical data visualisation; the QGIS was applied for geological mapping. We analyzed 2000 samples of the earthquake events obtained from the IRIS seismic database with a 25-year time span (1997--2021) in order to map the seismicity. The approach to linking geological, topographic, and geophysical data using GMT scripts aimed to map correlations among the geophysical phenomena, tectonic processes, geological setting, seismicity, and earthquakes. The practical application of the GMT scripts consists in automated mapping for the visualization of geological risks and hazards in the mountainous region of the Venezuelan Andes. The proposed method integrates the approach of GMT scripts with state-of-the-art GIS techniques, which demonstrated its effectiveness as a tool for mapping spatial datasets and rapid data processing in an iterative regime. In this context, using GMT and GIS to find similarities between the regional earthquake distribution and the geological and topographic setting is essential for hazard risk assessment. This study can serve as a basis for predictive seismic analysis in geologically vulnerable regions of Venezuela. In addition to a technical demonstration of GMT algorithms, this study also contributes to geological { and }geophysical mapping and seismic hazard assessments in South America. We present the full scripts used for mapping in a GitHub repository.}

\keyword{risk; hazard; sustainability; earthquake; seismicity; cartography; GMT; geophysics} 

\PACS{91.10.Da; %MDPI: Please check if this is necessary, if it is, please give an acceptable reason.  If not, please remove them. <-- Answer: yes, the PACS, MSC and JEL codes are necessary for indexation of the article and classification of topics in the libraries.
 91.10.-v; 91.10.Jf; 91.10.Op; 91.30.Dk; 93.85.Pq; 91.70.-c}
\MSC{86Axx; %MDPI: Please check if this is necessary, if it is, please give an acceptable reason.  If not, please remove them. <-- Answer: yes, codes are necessary, answered above.
 86A04; 86A15; 86A30; 86A60; 86-XX; 86-04; 86-08}
\JEL{Y91; %MDPI: Please check if this is necessary, if it is, please give an acceptable reason.  If not, please remove them. <-- Answer: yes, codes are necessary, answered above.
 Q00; Q01; Q2; Q20; Q24; Q3; Q35; Q5; Q50; Q51; Q54; Q55; Q56; C6; C61; C63}

%%%%%%%%%%%%%%%%%%%%%%%%%%%%%%%%%%%%%%%%%
\begin{document}

\section{Introduction}
\unskip

\subsection{Background and~Motivation}

\textls[-15]{Nowadays, geological hazards are acknowledged as some of the most disastrous and high-risk events in mountainous regions. These involve a series of negative consequences for nature (landslides, tsunami, rock falls, land mass movements) and complex social problems (destruction of houses and infrastructure, injuries and losses of human lives) caused by seismic risks \citep{Laffaille,Schmitz2020,Weber}. Early prediction of seismic risks in geologically vulnerable regions has become more valuable in a large variety of time-critical geospatial applications focused on evaluating  geological hazards {for the mitigation of }earthquake hazards at the regional scale. The~integration of geological and geophysical data is an effective {approach for }risk assessment in mountainous regions, for the prevention of damages and danger, and for serving for early warning, long-term monitoring{, management, }and mitigation of seismic~hazards. }

{Recently, }advances in cartographic development and mapping have resulted in the realization of an important application of spatial data analysis: the prediction and prognosis of seismic activities or imminent earthquake events from   events observed in the past~\mbox{\cite{KHAN2020101642,MESGAR2017454,NATH2005329}}. {With enough data on earthquake events, these methods }enable{ an accurate }prognosis of seismicity with the aim of  mitigating hazard risks{ }\cite{,KWAG20227230,KWAG2021107809}. Early detection of seismic risks can be solved by evaluating and visualizing different geophysical determinants that affect disaster fatalities at the country level, namely, seismic hazard strength, population density, distribution of locations{, }demographic and socioeconomic data, ands geological background layers~\cite{LIU20111084,ijerph191912351,rs14174269,ijerph191811469}. For~instance, a~systematic examination of the correlations and links among various geophysical and geological determinants underlying disasters {is useful for }the prevention of{ }socioeconomic risks in{ }geologically unstable zones~\cite{YOUSUF2020100064,KAMATA2022103253}. In earthquake analysis, the~capability of predicting {earthquake} events, their magnitude, and their focal depth is highly desirable for real-time risk assessment~\cite{SHAO2022100177,YAGHMAEISABEGH2022107173}{.} 

To build{ }effective geological{ }modeling{ approaches }for seismic hazard prediction,{ }real-time mapping applications{ are beneficial for earthquake} prognoses{. }However, obtaining an {automated approach to }cartographic dataset {processing }for seismic prediction{ is usually difficult, resulting in biases toward a particular GIS software. }The approach presented in this study aims to solve {this }problem{ }by introducing an integrative study oof scripts developed in the Generic Mapping Tools (GMT) scripting toolset~\cite{Wessel},{ in addition to the traditional mapping}. This framework, which is different from the {conventional state-of-the-art } GIS {methods}, presents a {rapid workflow of} mapping {based on the data processing from the} console using{ }command line and codes written in the GMT syntax. Although~{advanced methods of }geophysical mapping  of such seismically active regions as South America {present }a very important task, {the scripting of cartography remains }a  new topic. To~the best of our knowledge, the~majority of similar studies used the conventional GIS rather than scripts for the visualization of {geophysical }data~\cite{Castillo}{.}

While conventional methods of GIS enable the plotting the maps in a traditional regime, this often requires manual operations and is not suitable for automated mapping. In~such a way, it limits the types of cartographic activities that can be used for operative geological analyses and seismic prediction. In addition, data handling is limited by {the }restricted compatibility of formats accepted in GIS{. }In contrast, scripts present a rapid and repeatable procedure of automated mapping that can be integrated {into }real-time risk assessment and hazard analysis. {As a response to these needs, }GMT presents a more flexible {cartographic }approach and{ }versatile strategy of {spatial }data processing, and it is aimed at {seismic }risk assessment{, geological exploration, and }management by using multi-source~datasets. 

\subsection{Contemporaneity and~Objectives}

Integrated mapping using various multi-source data has the objective of visualizing a complex interplay among the geological, topographic, tectonic, and geophysical{ parameters. Regional effects of these settings cause }the Caribbean region and{ }Venezuela to be exposed to a high risk of earthquakes,{ }volcanos, and associated tsunami hazards. For~instance, the tectonic structures{ of the }deep basins and erosional troughs affect the submarine geomorphology,{ }which can be interpreted using geophysical and gravimetric measurements{. }Earthquake-related aftershocks, hazards, and{ }landslides {result in  damage to the }geotechnical infrastructure and bridges~\cite{Salcedo}, demolished buildings{, }loss of human lives~\cite{Rodriguez2018}, and aseismic slips and creeps~\cite{Pousse2016}{. Such consequences }affect both nature and the social {sectors }of Venezuela and require regular monitoring {and} mapping{ for }risk and vulnerability~assessment. 

{The }visualization of the information on these processes is {important }for seismic prognosis and earthquake modeling. Recent research demonstrated that geological and geophysical mapping creates a basis for natural hazard risk assessment~\cite{su10020304,su131810232,Lemenkova202212,app11219919,ijgi10030118}. Therefore, the~data derived from hazard risk mapping can be used as background information and for the recognition of complex geological actions and processes{. }Previous cartographic investigations were mostly  carried out by using traditional GIS in geological and environmental studies \citep{Gonzalez,Hackley,Schmitz2021,Sun}. However, none of them focused on an integrated approach {that would combine }various cartographic techniques for the geophysical mapping of~Venezuela. 

{Scripting techniques applied to }geological mapping have not been adequately presented in the existing literature{ in comparison with the traditional GIS}. At~the same time, scripts{ }employed for visualization of  complex geophysical and geological phenomena enable automation{ }and increase the precision of {the processing of }cartographic data. The~use of machine-based methods enables one to accurately map and highlight{ }correlations between geological and geophysical phenomena{. }With this aim,~GMT {scripts }in combination with QGIS and its plugins{ }were applied for the geophysical mapping of Venezuela {in order to visualize the }regional seismicity {in connection }with the geological structure {and geophysical setting}. 

{The }cartographic objective {of this study }is to demonstrate the {functionality }of GMT scripts in  combination with GIS techniques. The~geological aim is to highlight the correlations between the geological setting, topography,  geophysical anomalies, and regional seismicity through the analysis of the distribution, magnitude, and depth of earthquakes. {To this end, this }study applied  GMT  cartographic scripting techniques and QGIS {in order to process }high-resolution open-source raster grids and vector data. We used a variety of datasets to build the project with observed geological and geophysical processes and structures{. The~primary }context {of this work is }the analysis of seismic events {by using }the visualization of earthquake locations{ }in the~region. 

\section{Study~Area}
This study focuses on Venezuela{, }one the of the most seismically active regions{ }of {South America. The~spatial extent of the study area was} 59.%MDPI: We changed comma to dot. Please confirm this revision. <-- Answer: yes, we confirm.
5°--74° W, 0°--12.5° N{ }(Figure \ref{fig1}). 

\vspace{-10pt} {}
\begin{figure}[H]
	\hspace{-10pt} {}\includegraphics[width=11.0 cm]{Fig_01}
	\caption{Topographic map of the region of Venezuela. Mapping: GMT. %MDPI: Please ensure that permission has been obtained and there is no copyright issue. <-- Answer: yes, we confirm that the maps are made by the authors and there is no copyright issue.
 Source: authors. \label{fig1}}
\end{figure}  

{Its }location in the Andes{ exposes Venezuela }to a high risk of seismic hazards.
Geologically, the~most important factors that have a significant influence on the seismicity in Venezuela and its surroundings include lithospheric plate subduction and the associated complex geological processes{ of the }oblique collision of the Caribbean arc and a zone of active deformation in the eastern offshore Trinidad area~\cite{ESCALONA20118}. This is reflected in the formation of the Venezuelan Andes as a N50°E-oriented mountain belt extending from the Colombian border in the SW to the Caribbean Sea in the NE with the associated Boconó and Valera thrust faults~\cite{DHONT2012245}. Seismotectonic activities in Venezuela can be characterized by a spatio-temporal composition of the constituent geophysical and geological forces{, their }evolutional processes, and{ their }complex interactions. Such linkages are related to the actual location and distribution of seismic sources of the earthquakes in the Andean region with~regard to the{ }tectonics and overall dynamics of the Earth's crust in South~America. 

\textls[-15]{The {high }frequency and magnitude of earthquakes in Venezuela are associated with {complex }geological and geophysical settings{; Figure} \ref{fig2}. {In addition, the }tectonic interactions of the lithospheric plates{ trigger }seismic activity and tectonic-related hazards {that }affect the lives of people and infrastructures \citep{Masy,Perez2018,Suarez}.{ }The Caribbean--South American plate boundary is a very wide tectonically active zone~\cite{Audemard1993,Backe} with {recorded }seismicity at various depths. {Thus}, anomalous series of earthquakes took place on the east coast of  Trinidad with a reported depth of the main shock of 53 km~\cite{Marshall}. A seismicity deeper than 40 km was reported on the southern end of the Lesser Antilles subduction zone and{ }the active island arc~\cite{Reinoza}. }

\vspace{-2pt} {}
\begin{figure}[H]
	\hspace{-3pt} {}\includegraphics[width=13.5 cm]{Fig_02}
	\caption{\hl{Geological map} %MDPI: The contents of this figure are not legible. In order to convert a clear PDF document, whilst retaining its high quality, we kindly request the provision of figures and schemes at a sufficiently high resolution (min. 1000 pixels width/height, or a resolution of 300 dpi or higher).
 of the region of Venezuela{. }Mapping: QGIS. Source: authors.\label{fig2}}
\end{figure}  

{Venezuela} is one of the most  countries that is most vulnerable to earthquakes{ in the Amazon Basin}. {The }geological hazards {in this region }are exacerbated by social vulnerability due to the anthropogenic exposure, with{ a }frequently injured and affected population {living }in {the }high mountainous regions of Andes{, which }increases the population who is at risk during earthquakes{. }Furthermore, catastrophic earthquakes produce a chain effect that triggers tsunamis, which amplifies the damage to the environment and human society. For~example, it is reported that since the 19th century{, }tsunamis in the Caribbean Sea have resulted in the loss of several thousands of lives  and the exposure of many people  to risk and damage from  landslides~\cite{Harbitz}. 

Venezuela is{ }located in the north of the South American continent, where a collision of the South American and Caribbean tectonic plates{ causes }complex regional geological patterns{ }(Figure \ref{fig2}). The~geodynamics of the interacting plates result in the formation of {}thrust belts and foreland basins in northern Venezuela that are {shaped }by the interplay of the geological layers.{ }Outcrops of {}Precambrian (pC), Quaternary (Q), Triassic (T), and Paleozoic--Mesozoic intrusives (PMi)  dominate in the central regions of the country{, }while regional outcrops of {}Cretaceous (K) and Paleozoic metamorphics (PZm) and Mesozoic--Cenozoic intrusives (MCi) are {found }in the Falcón Basin (Figure \ref{fig2}). 

\textls[-25]{The northwestern and northeastern regions{ of Venezuela are }prone to earthquakes~\cite{Audemard,Perez}} due to the tectonic activity of the lithospheric plates \citep{Barr,Hayes}. {According to the Incorporated Research Institutions for Seismology (IRIS) database, the~}repetitive occurrence of seismic events, such as volcanism and earthquakes, in northwestern{ }Venezuela {and} Trinidad and Tobago was recorded as {a }series of{ }events {over the }last 25 years. The~ongoing effects of earthquakes in Venezuela \citep{Audemard2006,Baumbach,Colon,Pousse2018,Rial,Yamazaki} and coastal regions of Trinidad and Tobago \citep{Deville,Latchman} have made geophysical mapping an important task {for} seismic hazard mapping, earthquake {prognosis, }and risk assessment~\cite{Bucci}.


The region of the  Caribbean Sea  is notable for its complex geologic setting (\mbox{Figure \ref{fig3}}), where several tectonic plates experience collision, interaction, and subduction---these are the Caribbean, Cocos, Nazca, and South American plates. This results in {a }high seismic activity in the region{. The} influence of tectonic processes{ is }reflected in {a variety of topographic forms, geological structures, and geomorphic settings} of{ }Venezuela{, }as{ }discussed {earlier } \citep{Diebold,Fajardo,Lemenkova2020b,Perez1981}. The~geodynamic heterogeneity of the upper mantle has resulted in a specific mountain geomorphology with notable features, such as the extent of the Tobago Trough, the Lara--Falcón Basin, and the crustal structure of the Cordillera de Mérida in southwestern Venezuela, which are~important segments of {the }Andean orogeny~\cite{AVILAGARCIA2022103853}. 
The distribution of earthquakes in this region is strongly controlled by the tectonic and geophysical setting, i.e.,~closeness to the margins of the colliding South American and Caribbean tectonic plates. Other{ important factors include} the mountain geomorphology and topographic variability, which are shaped by the geophysical--geological forces that affect the variations in depth and magnitude of seismic~events. 

\begin{figure}[H]\vspace{-6pt}%\centering
	\includegraphics[width=13 cm]{Fig_03}
	\caption{Geological provinces in  the region of Venezuela{. }Mapping: QGIS. Source: authors.\label{fig3}}
\end{figure} 



{Active} tectonic movements cause the formation of  orogenic mountain belts and topographies. {In addition, regional }bathymetry {is }complicated {by the }deep-sea trenches and troughs in the margins of the Pacific Ocean and the Caribbean Sea along the subduction path \citep{Gough,Lemenkova2020a}. The complexity of the  geological history is especially notable in  NW Venezuela---in the Falcón Basin---where{ }interactions between the lithospheric plates and subducting Caribbean slab cause  associated seismic activity and geological processes, such as {an }outcrop of {}igneous intrusive bodies, crustal thinning, decreased Moho depth, gravity anomalies, {and} rifting~\cite{Bezada}. {Previous studies }\cite{Mazuera} identified crustal thinning beneath the Falcón Basin along the western extension of the Oca--Ancón fault system{, }which was interpreted as a back-arc basin, as well as~suture zones between the Proterozoic and Paleozoic provinces (Ouachita--Marathon-related suture) and~Paleozoic and Meso-Cenozoic terranes (peri-Caribbean suture), where the seismic \mbox{velocity changes{.} }

%%%%%%%%%%%%%%%%%%%%%%%%%%%%%%%%%%%%%%%%%%
\section{Methodology}
\unskip

\subsection{Data}
Topographic, geophysical, seismic, and geological data{ }were taken {from  open sources }in order to present a multi-disciplinary cartographic analysis of Venezuela. The~topographic data were collected {from }the open-access repository {of the }General Bathymetric Chart of the Oceans (GEBCO) {in }NetCDF, and they were collected {in }raster {format with }grid {cells }\cite{GEBCO}{. The~GEBCO }presents the most comprehensive data on the topography and bathymetry of the Earth {with a 15 arc second resolution } \cite{Schenke}.~GEBCO data{ are }widely used in the geosciences due to their {robustness, i.e., their high }resolution{, }precision,{ and reliability as a cartographic data source }\citep{Mcmichael,Cortina}. 

Geological data (Figures \ref{fig2} and \ref{fig3}) were {collected from the United States Geological Survey (USGS). The~geological maps were processed as .shp vector layers} by using{ }the {open-source QGIS software due to its} compatibility{ }with the ArcGIS data format. The~visualization of these maps followed the traditional mapping process of  QGIS~\cite{Lemenkova2022a}. In~the first stage, vector layers of the geological provinces and outcrops were visualized within a graphical user interface (GUI). In~the second stage, the~study area was designated by using an overlay of the Digital Chart of the World (DCW), which is a comprehensive 1:1,000,000 scale vector base map uploaded into the QGIS~project. 

\textls[-5]{{The geoid map was visualized} based on {data }collected from  EGM2008~\cite{Pavlis}. {The free-air gravity maps and vertical gradient were based on  open-source datasets available from the public repositories of the Scripps Institution of Oceanography and described in} \cite{Sandwell}. {The seismic dataset used to plot maps of earthquakes was obtained from the Incorporated Research Institutions for Seismology (IRIS) database by using the spatial extent of Venezuela and the neighboring areas of Trinidad, Tobago, and the surroundings. This~dataset covered the earthquake events in the region of Venezuela   for the period from 1997 to 2021 and included earthquakes with magnitudes from 1.9} to 7.3 $M_L$ on the {Richter scale.}}

\subsection{Methods}
Maps {showing the }topographic, seismic, and geophysical {setting of Venezuela }\linebreak{}(Figures \ref{fig1} and \ref{fig4}--\ref{fig7}) were made by using the scripting techniques{ }of{ }GMT~\cite{Wessel} by following the existing {methodology } described by \citep{Lemenkova2022b,Lemenkova202217}. The~methodological workflow included {five }algorithms in GMT codes, {which are presented in  Appendix \ref{appendix}  of this article.} Prior to data visualization and mapping, the{ scripts }were written using {GMT} syntax and {organized }in the Xcode environment. {As mentioned in the previous subsection, we processed the high-resolution data by using GMT. Conventionally, the~}language of GMT was used {in several modules for computing and processing raster data }in order to plot{ }cartographic elements. {The} area of interest was first selected and cut from the GEBCO grid by using the following {snippet of }code: $gmt grdcut GEBCO\_2019.nc -R286/300.5/0/12.5-Gve\_relief.nc$. %MDPI: Please confirm if the italics should be retained. Please check the whole paper. <-- Answer: yes, the italics separates the fragments of the programming code from the text; in this way it helps to improve the readability.

After the first step, all other  substantial cartographic elements{ }were added into a script by using the following GMT modules: psbasemap, grdimage, psscale, pstext, psclip, and grdcontour. {Specifically, these included} images that were visualized in selected color palettes, grids, clipped insertions of regions into a global map, % Please confirm that meaning has been retained. <-- Answer: yes, correct, we confirm.
 coastlines, borders, rivers, bar legends {on the maps}, scale bars{, }text annotations, time stamps{, etc. }The technique of scripting was applied by using  existing detailed descriptions~\cite{Lemenkova2021b}. For~example, the~'grdcontour' module was used to compute, model, and visualize topographic isolines by using quantitative DEM data as follows: `$gmt grdcontour ve\_relief1.nc -R -J -C1000 -Wthinner,darkbrown -O -K >> \$ps$'. The~purpose of the geoid visualization (Figure~\ref{fig4}) in relation to the geophysical differentiation was to investigate how the gravity values responded to the effects of topographic elevations and bathymetric~depressions. 

\textls[-15]{While GEBCO data were gathered in NetCDF format ($GEBCO\_2019.nc$), the~geoid data first needed to be converted  from the .adf format{. This }was done with the following code of GMT: `$gmt grdconvert n00w90/w001001.adf geoid\_02.grd$'. The  extent and range of the data were examined through the {Geospatial Data Abstraction Library }(GDAL) and interpreted for the color palette as follows: `$gmt makecpt -C33_blue\_red.cpt -T-55/25 > colors.cpt$'.} The geoid was modeled by using the following code: `$gmt grdimage geoid\_02.grd -Ccolors.cpt -R286/300.5/0/12.5 -JM6.5i -P -Xc -I+a15+ne0.75 -K > \$ps$'. {The visualized data are shown in Figure} \ref{fig4}. 
\vspace{-3pt} {}
\begin{figure}[H]
	\hspace{-10pt} {}\includegraphics[width=10 cm]{Fig_04}
	\caption{\hl{Geoid model} %MDPI: The contents of this figure are not legible. In order to convert a clear PDF document, whilst retaining its high quality, we kindly request the provision of figures and schemes at a sufficiently high resolution (min. 1000 pixels width/height, or a resolution of 300 dpi or higher).
 of the region of Venezuela{. }Mapping: GMT. Source: authors.\label{fig4}}
\end{figure} 

The gravity grids and {the }geophysical data (Figure \ref{fig5} and \ref{fig6}) were also converted from the initial raw IMG format into the GRD format, {which is compatible} with  GMT{. The~code was executed }as follows: `$gmt img2grd grav\_27.1.img -R286/300.5/0/12.5 -Ggrav\_VE.grd -T1 -I1 -E -S0.1 -V$' (for {Figure }\ref{fig5}). The~same procedure was repeated for{ Figure }\ref{fig6} (vertical gradient of gravity){ by following }the examples{ described in  existing }work~\cite{Lemenkova2021d}. Annotations were added to all of the maps by using the 'pstext' module of GMT, as in the following example: `$gmt pstext -R -J -N -O -K -F+f10p,0,black+jLB+a-0 >> \$ps << EOF 294.20 8.24 Ciudad Bolívar EOF$'{. }The data were converted into {the }‘ngdc’ format, visualized with the 'psxy' module, {and mapped, as shown in Figure} \ref{fig5}. 

The earthquakes {in }Figure \ref{fig7} (map of seismicity) were plotted by using the `psxy' module of GMT with the following code: `$gmt psxy -R -J quakes\_VE.ngdc -Wfaint -i4,3,6,6s0.05 -h3 -Scc -Csteps.cpt -O -K >> \$ps$'. Here,  ‘i4,3,6,6s0.05’ indicates the numbers of the columns in the table ($quakes\_VE.ngdc$) that were converted from the original CSV from  IRIS. The~volcanoes were visualized with the following code: `$gmt psxy -R -J volcanoes.gmt -St0.4c -Gred -Wthinnest -O -K >> \$ps$’. The~complex legend showing the magnitude of events was plotted by using the following code: `$gmt pslegend -R -J -Dx1.5/-3.0+w17.8c+o-2.0/0.1c  -F+$\emph{pthin+ithinner+gwhite -O -K << FIN >> \$ps H 10 Helvetica Seismicity: earthquakes magnitude (M) from 1.9 to 7.3 N 9 S 0.3c c 0.3c red 0.01c 0.5c M (7.1-7.3) <...> FIN}'.

\vspace{-3pt} {}
\begin{figure}[H]
	\hspace{-16pt} {}\includegraphics[width=10.5 cm]{Fig_05}
	\caption{\hl{Free-air} %MDPI: The contents of this figure are not legible. In order to convert a clear PDF document, whilst retaining its high quality, we kindly request the provision of figures and schemes at a sufficiently high resolution (min. 1000 pixels width/height, or a resolution of 300 dpi or higher).
 gravity model of the region of Venezuela{. }Mapping: GMT. Source: authors.\label{fig5}}
\end{figure} 
\vspace{-12pt} {}

\begin{figure}[H]
	\hspace{-12pt} {}\includegraphics[width=10.5cm]{Fig_06}
	\caption{\hl{Vertical} %MDPI: The contents of this figure are not legible. In order to convert a clear PDF document, whilst retaining its high quality, we kindly request the provision of figures and schemes at a sufficiently high resolution (min. 1000 pixels width/height, or a resolution of 300 dpi or higher).
 gradient of gravity in the region of Venezuela{. }Mapping: GMT. Source: authors.\label{fig6}}
\end{figure}  


\vspace{-6pt} {}
\begin{figure}[H]
\hspace{-12pt} {}	\includegraphics[width=12.0 cm]{Fig_07}
	\caption{Seismic map of earthquakes in the region of Venezuela{. }Mapping: GMT. Source: authors.\label{fig7}}
\end{figure}   

{In practice, we fused the data from seismic events to the topographic grid to show the locations of the events on the map with their attributes, which showed magnitudes ranging from 1.9 to 7.3. In~our maps, we demonstrated that the density of the earthquake events visibly increased along the lines of the boundaries of lithospheric plates  (thick purple line on the map in Figure } \ref{fig7}){, with a maximal concentration of  events in the Andes (Cordillera de Merida) and the region of Trinidad and Tobago}. 

The full algorithms of the GMT codes are available in the authors' GitHub repository:\\ \href{https://github.com/paulinelemenkova/Mapping_Venezuela_GMT_Scripts}{https://github.com/paulinelemenkova/Mapping\_Venezuela\_GMT\_Scripts}, accessed on 1 November 2022. %MDPI: Please add the access date (Format: Date Month Year). e.g., (accessed on 1 January 2020). Please note that the access date should be earlier than the submission date 4 November 2022. <-- Answer: added the access date (1 November 2022).


%%%%%%%%%%%%%%%%%%%%%%%%%%%%%%%%%%%%%%%%%%
\section{Results}

Figure~\ref{fig1} was plotted by using the GEBCO raster grid, which was selected due to its quality and reliability.{ A }precise topographic{ map is essential to a fundamental }understanding of the geophysics and tectonics {of a region. Topographic maps enable the analysis of earthquake locations and the evaluation of the }associated effects on ocean and terrestrial geomorphology. {Therefore, the~analysis of  topography is necessary for} the mapping of seismicity, risk assessment, and tsunami wave propagation. Here, the topographic and bathymetric elevations of the study area encompassing Venezuela and the islands of Trinidad and Tobago  ranged from $-$4948 to 5536 m{ }according to the GEBCO~grid. 



Figures~\ref{fig2} and \ref{fig3} indicate the spatial extent of the geological provinces and outcrops in Venezuela. The~map in Figure~\ref{fig2} shows that most of the territory of Venezuela is covered by {the }Guyana Shield (aquamarine color) in the southeastern parts of the country. The~Barinas--Apure sedimentary basin is located in{ }Llanos{ }and the region of the Andean foothills{ in }western Venezuela~\cite{Callejon}. The~Barinas--Apure basin (yellow color on the map) includes oil fields{, which comprise }the third most important  area of production in Venezuela after the Maracaibo and Eastern Venezuela basins{ }according to the petroleum accumulation. The~Falcón sedimentary basin, which produces indigenous Miocene oil from structural and stratigraphic accumulations~\cite{Rodriguez}{, }is colored in slate blue on the~map. 

The South Caribbean deformed belt (purple color in Figure~\ref{fig2}) represents a submarine prism formed between the subducting tectonic plates in the Colombian and Venezuelan basins and the arc terranes along the northern rim of South America~\cite{Mann}. Its {southwestern} part, the~Sinu fold belt{, }forms an accretionary wedge in the {area }of{ }subduction of the Caribbean Plate under {the }South American Plate~\cite{Rodriguez2021}. Other important geological provinces include the Guyana Basin, which stretches along the passive margin of {northeastern} South America, and the Maracaibo Basin (green, Figure~\ref{fig2}). The~Cariaco Basin (beige color in Figure~\ref{fig2}) is an East--West-stretching basin~\cite{Schubert} situated in the Gulf of Cariaco, which is {located }on the continental shelf of the Caribbean Sea, off~the eastern coast {of the surroundings of }Barcelona,{ }Figure \ref{fig1}). 

The Tobago Trough (chartreuse green color in Figure~\ref{fig2}) is a modern marine forearc basin of the Lesser Antilles arc and is located above the sedimentary succession{. It is }at least 10 km thick and is a probable oceanic basement{ that was }formed by sediments from various stratigraphic sequences---Quaternary, early Pleistocene--late Miocene, Miocene, and late Eocene{. In addition, it }may also include mid-Cretaceous or older {intrusive suites }\cite{Speed}. The~{distribution of the }Quaternary sediments covering  most of the northwest of the country (dark pink color in Figure~\ref{fig3}) corresponds{ }with  previous studies on Plio-Quaternary extension in the Venezuelan Andes{. Thus, this }shows{ }the extensional structures corresponding to  elongated tilted blocks, with the geometry and kinematics of the structures corresponding to the earlier syn-orogenic extension~\cite{Dhont}. 


Figure~\ref{fig4} shows the variations in the geoid heights over the study area. The~highest values for the geoid (up to 25 m, bright red {colors) }are visible in the southwestern region of the map {at the border with }Colombia{. However, }the majority of {the territory of }Venezuela covers the extent from $-$50 to 8 m, with the highest values being in the mountains{, showing a }clearly visible correlation with topographic elevations. The~undulations of the geoid in the Guyana Highlands correspond to the higher topographic elevations on the shield, with table-like  `tepuis' mountains. % Please confirm that meaning has been retained. <-- Answer: yes, we confirm, the meaning has been retained.
 In general, the~marine areas of the Caribbean Sea have  lower geoid heights (blue-colored areas in Figure~\ref{fig4}), which  correspond well with the bathymetric depressions reflected in the geophysical setting of the~Earth.

{The comparison of the }geoid mode in Figure~\ref{fig4} with the topographic and geological maps in Figures~\ref{fig1} and \ref{fig2}{ shows }the elevated values of the geoid{ that }are notable in the region of the Guyana Shield. The~Guyana Shield is a Precambrian geological formation on the NE coast of  South America and {is }an area of special importance in the historical geology of Venezuela due to {the }high potential for mineral exploration~\cite{Gibbs}. {Aside from the geological features, the~}environmental importance of this unique region consists in the highest biodiversity in the world,{ as it includes the numerous }endemic species~\cite{Hollowell} {of }the tropical forests of Amazonia~\cite{Hammond}. 

Figure~\ref{fig5} shows the variations in the gravity over the study area. The~notable minimum of the free-air gravity  coincides % Please confirm that meaning has been retained. <-- Answer: yes, we confirm, correct.
 with distribution of the Guyana Basin, the {southeastern} part of  Venezuela,{ and }Trinidad and Tobago{. In~particular, this }points to the tectonic basement depression formed during {the }Jurassic rifting, {as well as }Lake Maracaibo, an~important geological region of Venezuela. The~{northeastern }coasts of  Lake Maracaibo are occupied by the Bolivar Coastal Field, the~largest oil field in South America{. With~a basin area of 50,000~km$^{2}$ and oil production of $30 \times 109$ bbl %MDPI: Please check if this should be a multiplication sign. <-- Answer: yes, correct, changed to \times
, it is }the second{ }most prolific hydrocarbon basin in the world after the Middle Eastern{ hydrocarbon basins }\cite{Escalona}.   

 Maracaibo Lake has a wide variety of oil reservoir rocks with structural traps in the tectonic units, e.g.{, }normal or inverted faults on the South American plate~\cite{Bockmeulen}. The~oil formed in the Maracaibo Basin is mostly heavy black oil with asphaltene deposition~\cite{Escobar,Memon}. The~gravity in  Lake Maracaibo reaches below $-$200 mGal in correlation with  the bathymetric depression (Figure \ref{fig1}) and the geological extent of the Maracaibo Basin (Figure \ref{fig2}).{ The }highest values for gravity reach +200 mGal and  correspond well with the topographic elevation of the Cordillera de Mérida, Columbian Andes, Guiana Highlands, and Sierra Parima. The~Amazon and Orinoco Flow Basins mostly have gentle fluctuations in gravity, with values in the range between $-$25 and 30 mGal (Figure \ref{fig5}).

Figure~\ref{fig6} refers to a vertical{ }gradient of normal gravity as an approximation in which a gravity measurement point, which is almost never placed on the reference ellipsoid’s surface{, }is adjusted according to geophysical corrections \citep{Li,Yilmaz}. The~details regarding the distribution of vertical gravity values over{ }Venezuela,  Trinidad, and Tobago showed that the lowest values (below $-$100 mGal) mostly correspond to the extreme amplitude of topographic changes{, e.g.,} the abrupt borders between the mountains/hills and the margins of the bathymetric depression{, }while the highest values (over 100 mGal) were detected on the tops of the mountains, {which }corresponded well{ }with the local geomorphology of the{ }region.

The distribution of the earthquakes and the seismicity of  Venezuela,  Trinidad, and Tobago (Figure \ref{fig7}) depend on the proximity of the area to a border of  lithospheric plates  with active margins{, }such as the South American and Caribbean plates{, }which are depicted in the map with purple thick lines{. }The second factor includes   active geological processes (faults) and~the geomorphology. Hence, the Cordillera de Mérida had the majority of the detected earthquakes after the active zone of  Trinidad and Tobago. {Most }of the earthquakes{ mapped in }Venezuela {belonged} to the  ‘shallow’ category, that is, they had a depth between 0 and \mbox{70 km.} Earthquakes at depths from 74 to 155 km were located offshore of  Sucre, Venezuela, and those that were detected in  Trinidad and Tobago were predominantly shallow{, }not exceeding 70 km. The~deepest earthquakes in the inspected database were located in  northern Columbia (182--200 km).

\section{Conclusions}

In this paper, we  proposed two algorithms for the geophysical mapping of Venezuela, which included {}GMT-based scripts and traditional plotting by using QGIS, with the aim of analyzing the geohazard risk  of Venezuela.  Investigations and risk assessments of geohazards require advanced methods of mapping that are based on the programming of algorithms, as~demonstrated in this paper. The~operative use of correctly and effectively plotted maps, along with tabular data, supports reasonable and substantiated % Please confirm that meaning has been retained. <-- Answer: yes, correct.
 recommendations aimed at {the }prevention and mitigation of earthquake hazard risks in seismically active regions, such as the South American Andes. The~algorithms based on GMT enabled improved performance of the cartographic workflow {in comparison with that of the }conventional approaches{.} 

We demonstrated the technical advantages of {the }approaches to programming  in cartography, which included{ }scripting, as this is a~rapidly developing branch of geo-information. % Please confirm that meaning has been retained. <-- Answer: yes, correct. 
 Due to the high degree of automation, scripting is the best technical cartographic solution for seismic mapping, as~it enables rapid machine-based processing of data, which is {beneficial }for prognosis, predictive mapping, risk assessment, and hazard visualization in a real-time regime. Automation enables the rapid processing of  large datasets  and the effective use of  multi-source heterogeneous data. Aside from{ }that, the known advantages of scripting in cartography include machine-based graphics, which result in an aesthetically refined plotting by~GMT. 

Scripting cartographic methods for GMT were addressed to evaluate the geophysical and geological correspondences in a seismically active region of Venezuela and the Caribbean Sea basin. We used {}high-resolution geophysical, topographic, and geological datasets, such as GEBCO, EGM-2008 for geoid, gravity grids, and seismic data obtained from IRIS. The~results  corresponded with and supported those of previous studies on the geological and geophysical methods of gravity measurements, seismic observations, and earthquake monitoring in Venezuela and the surroundings \citep{Guenther,Schmitz2008}.  

This study presented a machine-based mapping workflow that contributed to the regional studies of the eastern Caribbean and Venezuela. The~GMT code snippets have been presented in a GitHub repository as~a demonstration of the technical possibilities of this toolset. The~results showed {}optimized {and improved }cartographic performance and demonstrated the advantages of the automating processing of geophysical datasets. Multi-source geophysical, geologic, seismic, and topographic data were used for a comparative analysis of the geophysical processes. The~tectonic evolution of the Caribbean Sea basin largely affected its geological setting, which could be tracked in the relevant geophysical data, indicating a high seismicity and  high earthquake~risk  in the Andes.

Since the eastern Caribbean is a tectonically complex area that includes a subduction zone, island arc, and active volcanoes~\cite{Ambeh}, a regional analysis and visualization of the varied settings were performed. Specifically, high seismicity was detected in the northwestern region of Venezuela and~the northeastern region of Trinidad and Tobago, which is consistent with the geological structure associated with the tectonic setting and~the related movements of the subduction of the Caribbean Plate. Although~continued research is needed in order to improve earthquake monitoring in the region of Venezuela, this study presents a contribution to the seismic study of the Caribbean and South America, as it shows the earthquakes that were detected according to the IRIS database based on a 25-year time span (1997--2021). 

A better understanding of the geological and tectonic factors affecting seismicity and regulating the appearance, magnitude, and focal depth of  earthquakes is critical for seismically active regions, such as the Caribbean and Venezuela. Practical applications of geological risk assessment in South America include the use of information by public administrations in city management and  urban planning for environmental management and the mitigation and management of hazard risk in regions prone to geological hazards, such as~Venezuela. 

{The main contributions that we have made and the contemporaneity % Please confirm that meaning has been retained. <-- Answer: yes, correct. 
 of the presented study are highlighted as follows:}

\begin{enumerate}
  \item Region: %MDPI: Please confirm if the bold should be retained. <-- Answer: bold deleted (it was not necessary)
 The region of Venezuela is at high risk of seismicity and earthquakes {because it is }located in the zone of the{ collision of the }South American and Caribbean tectonic plates and the Andean orogeny. This requires spatial analysis of the regions that are at risk.
  \item Geohazard: The risk assessment of geophysical hazards was based on{ the effective }integration and visualization of  multi-source data{, }which provided detailed insights into the environmental and geological settings of Venezuela and supported complex investigations in  seismically active regions of South America. 
    \item Data: Geological disasters are related to high seismicity, exposure to hazards, and~vulnerability of people at risk. This requires complex and detailed studies that summarize these factors and visualize regional geophysical settings{, as~shown in this paper}.
   \item Methods: The combination of  scripting for GMT and GIS is {a solid foundation} for cartographic data analysis. {An efficient method for processing }multi-format data  ensures hazard mapping and geophysical and geological visualization {with the aim of} preventing and evaluating  seismic hazards.
\end{enumerate}

{In this paper, we approached the problem of automated mapping in the geophysical sciences and the estimation of the  level of seismicity in the region of Venezuela for  risk assessment studies as a contribution from the new perspective of data analysis. The~scripting approach was mainly designed for the integration of data and the presented maps for studies of risk assessment,  though it is also valuable as a general-purpose mapping technique. The~limitations of  GMT may be constrained in  specific tasks, such as the analysis of remote sensing data  or image processing, where general-purpose languages, such as Python and R, \mbox{perform better.}}

{The  methodology that was presented and explained here for plotting maps with codes from a console can be potentially extended to other regions of the world, since the approach to data analysis remains identical. It is possible to reuse the presented codes for other seismically active regions of the world. The~only modifications of the method would include the spatial extent of the study area, that is, the cartographic coordinates for plotting the maps and for downloading the data from the IRIS catalogue. }

{The implementation of the GMT scripts for other regions implies that one  only needs to access datasets for a given region and to adjust the codes to this regional extent in order to map the data fairly well. We consider this cartographic scripting approach as a promising extension of our method for similar future  work on seismicity and risk assessment. We demonstrated that the performance and  visualization that were obtained are   effective and that the level of workflow automation is high. The~experimental results verify the usefulness of our approach for both of the following research objectives: the evaluation of the cartographic performance of GMT and the mapping of  seismicity in the region of Venezuela  and the northern Andes.} 

%%%%%%%%%%%%%%%%%%%%%%%%%%%%%%%%%%%%%%%%%%
\vspace{6pt} 

\authorcontributions{Supervision, conceptualization, methodology, software, resources, funding acquisition, and project administration, O.D.; writing---original draft preparation, methodology, software, data curation, visualization, formal analysis, validation, writing---review and editing, and investigation, P.L. All authors have read and agreed to the published version of the manuscript.}

\funding{This %MDPI: Please carefully check the funding information especially the grant number. <-- Answer: yes, funding information and grant number are correct.
 project was supported by the Federal Public Planning Service Science Policy or Belgian Science Policy Office, Federal Science Policy---BELSPO (B2/202/P2/SEISMOSTORM).} 

\institutionalreview{Not applicable.}

\informedconsent{Not applicable.}

\dataavailability{The GitHub repository contains the GMT scripts used for the mapping in this study: \href{https://github.com/paulinelemenkova/Mapping_Venezuela_GMT_Scripts}{https://github.com/paulinelemenkova/Mapping\_Venezuela\_GMT\_Scripts}, accessed on 1 November 2022. %MDPI: Please add the access date (Format: Date Month Year). e.g., (accessed on 1 January 2020). <-- Answer: added information on access date (1 November 2022).
} 

\acknowledgments{The authors thank the three anonymous reviewers for their careful reading, critical suggestions, and useful comments that helped improve an earlier version of this~manuscript.}

\conflictsofinterest{The authors declare no conflicts of~interest.} 
\vspace{-6pt} {}
\abbreviations{Abbreviations}{
The following abbreviations are used in this manuscript:\vspace{-6pt} {}\\

\noindent 
\begin{tabular}{@{}ll}
CSV & Comma-separated values \\
DCW & Digital Chart of the World \\
EGM & Earth Gravitational Models \\ 
GEBCO &  General Bathymetric Chart of the Oceans \\
GIS & Geographic Information System \\
GMT & Generic Mapping Tools \\
GUI & Graphical User Interface \\
QGIS & Quantum GIS \\
IRIS & Incorporated Research Institutions for Seismology \\
NetCDF & Network Common Data Form \\
NGDC & National Geophysical Data Center \\
USGS & United States Geological Survey \\
\end{tabular}}

%%%%%%%%%%%%%%%%%%%%%%%%%%%%%%%%%%%%%%%%%%
%% Optional
\appendixtitles{yes} % Leave argument "no" if all appendix headings stay EMPTY (then no dot is printed after "Appendix A"). If~the appendix sections contain a heading then change the argument to "yes".
\appendixstart
\appendix
\section[\appendixname~\thesection]{GMT Scripts}\label{appendix}
\subsection[\appendixname~\thesubsection]{GMT Script for Topographic Mapping}
\noindent\textbf{Listing A1:} GMT code used to plot Figure~\ref{fig1} (topographic map). 
\begin{lstlisting}[language=bash]
#!/bin/sh
# Purpose: shaded relief grid raster map from the GEBCO dataset (here: Venezuela)
# GMT modules: gmtset, gmtdefaults, grdcut, makecpt, grdimage, psscale, grdcontour, psbasemap, gmtlogo, psconvert
# http://soliton.vm.bytemark.co.uk/pub/cpt-city/arendal/tn/arctic.png.index.html
# GMT set up
gmt set FORMAT_GEO_MAP=dddF \
# Overwrite defaults of GMT
gmtdefaults -D > .gmtdefaults
#chsh -s /bin/bash
chsh -s /bin/zsh
gmt grdcut GEBCO_2019.nc -R286/300.5/0/12.5 -Gve_relief.nc
gmt grdcut ETOPO1_Ice_g_gmt4.grd -R286/300.5/0/12.5 -Gve_relief1.nc
gdalinfo ve_relief.nc -stats
# Minimum=-4947.963, Maximum=5535.625, Mean=162.220, StdDev=788.097
# create mask of vector layer from the DCW of country's polygon
gmt pscoast -R286/300.5/0/12.5 -Dh -M -EVE > ve.txt
# Make color palette
gmt makecpt -Carctic.cpt > pauline.cpt
# Generate a file
ps=Topography_VE.ps
# Make background transparent image
gmt grdimage ve_relief.nc -Cpauline.cpt -R286/300.5/0/12.5 -JM6i -P -I+a15+ne0.75 -t20 -Xc -K > $ps
# Add isolines
gmt grdcontour ve_relief1.nc -R -J -C500 -W0.1p -O -K >> $ps
# Add coastlines, borders, rivers
gmt pscoast -R -J -P \
    -Ia/thinner,blue -Na -N1/thickest,darkred -W0.1p -Df -O -K >> $ps
# CLIPPING
# 1. Start: clip the map by mask to only include country
gmt psclip -R286/300.5/0/12.5 -JM6.0i ve.txt -O -K >> $ps
# 2. create map within mask
# Add raster image
gmt grdimage ve_relief.nc -Cpauline.cpt -R286/300.5/0/12.5 -JM6.0i -I+a15+ne0.75 -Xc -P -O -K >> $ps
# Add isolines
gmt grdcontour ve_relief1.nc -R -J -C1000 -Wthinner,darkbrown -O -K >> $ps
# Add coastlines, borders, rivers
gmt pscoast -R -J \
    -Ia/thinner,blue -Na -N1/thicker,tomato -W0.1p -Df -O -K >> $ps
# 3: Undo the clipping
gmt psclip -C -O -K >> $ps
# Add color barlegend
gmt psscale -Dg286/-1.0+w15.2c/0.4c+h+o0.0/0i+ml -R -J -Cpauline.cpt \
    -Bg500a1000f100+l"Color scale 'arctic': global bathymetry/topography relief [-5000 to 4000, mixed, RGB, 110 segments]" \
    -I0.2 -By+lm -O -K >> $ps    
# Add grid
gmt psbasemap -R -J \
    --MAP_FRAME_AXES=WEsN \
    --FORMAT_GEO_MAP=ddd:mm:ssF \
    -Bpx4f2a2 -Bpyg4f2a2 -Bsxg4 -Bsyg2 \
    -B+t"Topographic map of Venezuela, Trinidad and Tobago" -O -K >> $ps    
# Add scalebar, directional rose
gmt psbasemap -R -J \
    -Lx12.7c/-2.2c+c50+w200k+l"Mercator projection. Scale: km"+f \
    -UBL/-5p/-65p -O -K >> $ps
# Texts
gmt pstext -R -J -N -O -K \
-F+jTL+f10p,21,darkred+jLB+a-0 -Gwhite@60 >> $ps << EOF
290.1 11.8 Peninsula
EOF
# insert map
# Countries codes: ISO 3166-1 alpha-2. Continent codes AF (Africa), AN (Antarctica), AS (Asia), EU (Europe), OC (Oceania), NA (North America), or~SA (South America). -EEU+ggrey
gmt psbasemap -R -J -O -K -DjBL+w3.3c+stmp >> $ps
read x0 y0 w h < tmp
gmt pscoast --MAP_GRID_PEN_PRIMARY=thinner,white -Rg -JG293/6N/$w -Da -Gbrown -A5000 -Bg -Wfaint -ESA+gpeachpuff -EVE+gyellow -Slightsteelblue3 -O -K -X$x0 -Y$y0 >> $ps
#gmt pscoast -Rg -JG12/5N/$w -Da -Gbrown -A5000 -Bg -Wfaint -ECM+gbisque -O -K -X$x0 -Y$y0 >> $ps
gmt psxy -R -J -O -K -T  -X-${x0} -Y-${y0} >> $ps
# Add GMT logo
gmt logo -Dx6.2/-2.9+o0.1i/0.1i+w2c -O -K >> $ps
# Add subtitle
gmt pstext -R0/10/0/15 -JX10/10 -X0.5c -Y5.0c -N -O \
    -F+f10p,Helvetica,black+jLB >> $ps << EOF
2.3 13.3 SRTM/GEBCO 15 arc sec resolution global terrain model grid
EOF
# Convert to image file using GhostScript
gmt psconvert Topography_VE.ps -A0.5c -E720 -Tj -Z
\end{lstlisting}

\subsection[\appendixname~\thesubsection]{GMT Script for Geoid Mapping}

\noindent\textbf{Listing A2:} GMT code used to plot Figure~\ref{fig4} (geoid modeling).

\begin{lstlisting}[language=bash]
#!/bin/sh
# Purpose: geoid of Venezuela
# GMT modules: gmtset, gmtdefaults, grdcut, makecpt, grdimage, psscale, grdcontour, psbasemap, gmtlogo, psconvert
# http://soliton.vm.bytemark.co.uk/pub/cpt-city/kst/tn/33_blue_red.png.index.html
# GMT set up
gmt set FORMAT_GEO_MAP=dddF \
# Overwrite defaults of GMT
gmtdefaults -D > .gmtdefaults
gmt grdconvert n00w90/w001001.adf geoid_02.grd
gdalinfo geoid_02.grd -stats
# Minimum=-70.856, Maximum=32.958, Mean=-24.768, StdDev=19.081
# Generate a color palette table from grid
gmt makecpt -C33_blue_red.cpt -T-55/25 > colors.cpt
# Generate a file
ps=Geoid_VE.ps
gmt grdimage geoid_02.grd -Ccolors.cpt -R286/300.5/0/12.5 -JM6.5i -P -Xc -I+a15+ne0.75 -K > $ps
# Add shorelines
gmt grdcontour geoid_02.grd -R -J -C1.0 -A2.0+f8p,0,black -Wthinner,dimgray -O -K >> $ps
# Add grid
gmt psbasemap -R -J \
    --MAP_FRAME_AXES=WEsN \
    --FORMAT_GEO_MAP=ddd:mm:ssF \
    -Bpx4f2a2 -Bpyg4f2a2 -Bsxg4 -Bsyg2 \
    --MAP_TITLE_OFFSET=0.8c \
    --FONT_ANNOT_PRIMARY=7p,0,black \
    --FONT_LABEL=7p,25,black \
    --FONT_TITLE=13p,25,black \
    -B+t"Geoid model (EGM-2008) of Venezuela and Trinidad and Tobago" -O -K >> $ps
# Add legend
gmt psscale -Dg286/-1.0+w16.5c/0.4c+h+o0.0/0i+ml+e -R -J -Ccolors.cpt \
    -Bg2f0.2a4+l"Color scale '33_blue_red': colour tables of IDL from the KDE scientific plotting tool [256, discrete, RGB, 252 segments, -T-55/25]" \
    -I0.2 -By+lm -O -K >> $ps
# Add scale, directional rose
gmt psbasemap -R -J \
    --FONT=7p,0,black \
    --FONT_ANNOT_PRIMARY=6p,0,black \
    --MAP_TITLE_OFFSET=0.1c \
    --MAP_ANNOT_OFFSET=0.1c \
    -Lx14.7c/-2.3c+c50+w200k+l"Mercator projection. Scale (km)"+f \
    -UBL/-5p/-65p -O -K >> $ps
# Add coastlines, borders, rivers
gmt pscoast -R -J -P -Ia/thinnest,blue -Na -N1/thick,brown -Wthick,darkslategray -Df -O -K >> $ps
# Texts
gmt pstext -R -J -N -O -K \
-F+jTL+f10p,21,darkred+jLB+a-0 >> $ps << EOF
290.0 10.5 (Coro)
EOF
# Add GMT logo
gmt logo -Dx7.0/-2.9+o0.1i/0.1i+w2c -O -K >> $ps
# Add subtitle
gmt pstext -R0/10/0/15 -JX10/10 -X0.1c -Y5.9c -N -O \
    -F+f10p,25,black+jLB >> $ps << EOF
3.0 13.6 World geoid image EGM2008 vertical datum 2.5 min resolution
EOF
# Convert to image file using GhostScript
gmt psconvert Geoid_VE.ps -A1.0c -E720 -Tj -Z
\end{lstlisting}


\subsection[\appendixname~\thesubsection]{GMT Script for Free-Air Gravity Mapping}
\noindent\textbf{Listing A3:} GMT code used to plot Figure~\ref{fig5} (free-air gravity modeling).
\begin{lstlisting}[language=bash]
#!/bin/sh
# Purpose: mapping free-air gravity anomaly of Venezuela
# GMT modules: gmtset, gmtdefaults, img2grd, makecpt, grdimage, psscale, grdcontour, psbasemap, gmtlogo, psconvert, pscoast
# http://soliton.vm.bytemark.co.uk/pub/cpt-city/pj/4/index.html
# GMT set up
gmt set FORMAT_GEO_MAP=dddF \
    MAP_FRAME_PEN=dimgray \
# Overwrite defaults of GMT
gmtdefaults -D > .gmtdefaults
gmt img2grd grav_27.1.img -R286/300.5/0/12.5 -Ggrav_VE.grd -T1 -I1 -E -S0.1 -V
gdalinfo grav_VE.grd -stats
# Minimum=-218.668, Maximum=942.457, Mean=5.694, StdDev=64.118
# Generate a color palette table from grid
gmt makecpt -Chaxby -T-200/200 > colors.cpt
# Generate a file
ps=Grav_VE.ps
gmt grdimage grav_VE.grd -Ccolors.cpt -R286/300.5/0/12.5 -JM6.0i -P -I+a15+ne0.75 -Xc -K > $ps
# Add isolines
gmt grdcontour grav_VE.grd -R -J -C50 -A100 -Wthinner -O -K >> $ps
# Add grid
gmt psbasemap -R -J \
    --MAP_FRAME_AXES=WEsN \
    --FORMAT_GEO_MAP=ddd:mm:ssF \
    -Bpx4f2a2 -Bpyg4f2a2 -Bsxg4 -Bsyg2 \
    -B+t"Free-air gravity anomaly for Venezuela, Trinidad and Tobago" -O -K >> $ps 
# Add legend
gmt psscale -Dg286/-1.0+w15.0c/0.4c+h+o0.0/0i+ml+e -R -J -Ccolors.cpt \
    --FONT_LABEL=7p,Helvetica,black \
    --FONT_ANNOT_PRIMARY=7p,0,black \
    --FONT_TITLE=8p,25,black \
    -Bg50f5a50+l"Color scale 'haxby' B. Haxby's color scheme for geoid & gravity [C=RGB -T-200/200]" \
    -I0.2 -By+l"mGal" -O -K >> $ps
# Add scale, directional rose
gmt psbasemap -R -J \
    -Lx14.0c/-2.3c+c50+w200k+l"Mercator projection. Scale (km)"+f \
    -UBL/-5p/-70p -O -K >> $ps
# Add coastlines, borders, rivers
gmt pscoast -R -J -P -Ia/thinnest,blue -Na -N1/thick,white -Wthin,darkslategray -Df -O -K >> $ps
# Texts
gmt pstext -R -J -N -O -K \
-F+jTL+f14p,32,darkgreen+jLB+a-330 >> $ps << EOF
290.8 7.8 Los Llanos
EOF
# Add GMT logo
gmt logo -Dx7.0/-2.9+o0.1i/0.1i+w2c -O -K >> $ps
# Add subtitle
gmt pstext -R0/10/0/15 -JX10/10 -X0.0c -Y4.8c -N -O \
    -F+f10p,25,black+jLB >> $ps << EOF
1.0 13.6 Global gravity grid from CryoSat-2 and Jason-1, 1 min resolution, SIO, NOAA, NGA.
EOF
# Convert to image file using GhostScript
gmt psconvert Grav_VE.ps -A0.5c -E720 -Tj -Z
\end{lstlisting}

\subsection[\appendixname~\thesubsection]{GMT Script for Vertical Gradient of Gravity}
\noindent\textbf{Listing A4:} GMT code used to plot Figure~\ref{fig5} (free-air gravity modeling).
\begin{lstlisting}[language=bash]
#!/bin/sh
# Purpose: geophysical mapping of Venezuela (vertical free-air gravity gradient)
# GMT modules: gmtset, gmtdefaults, img2grd, makecpt, grdimage, psscale, grdcontour, psbasemap, gmtlogo, psconvert, pscoast
# http://soliton.vm.bytemark.co.uk/pub/cpt-city/pj/4/index.html
# GMT set up
gmt set FORMAT_GEO_MAP=dddF \
# Overwrite defaults of GMT
gmtdefaults -D > .gmtdefaults
gmt img2grd curv_27.1.img -R286/300.5/0/12.5 -Ggrav_v_VE.grd -T1 -I1 -E -S0.1 -V
gdalinfo grav_v_VE.grd -stats
# Minimum=-488.036, Maximum=870.350, Mean=0.358, StdDev=43.643
# Generate a color palette table from grid
gmt makecpt -Chaxby -T-100/100 > colors.cpt
# Generate a file
ps=Grav_VE_v.ps
gmt grdimage grav_v_VE.grd -Ccolors.cpt -R286/300.5/0/12.5 -JM6.0i -P -I+a15+ne0.75 -Xc -K > $ps
# Add isolines
gmt grdcontour grav_v_VE.grd -R -J -C100 -Wthinnest -O -K >> $ps
# Add grid
gmt psbasemap -R -J \
    --MAP_FRAME_AXES=WEsN \
    --FORMAT_GEO_MAP=ddd:mm:ssF \
    -Bpx4f2a2 -Bpyg4f2a2 -Bsxg4 -Bsyg2 \
    -B+t"Vertical gravity gradient for Venezuela, Trinidad and Tobago" -O -K >> $ps  
# Add legend
gmt psscale -Dg286/-1.0+w15.0c/0.4c+h+o0.0/0i+ml+e -R -J -Ccolors.cpt \
    -Bg25f2.5a25+l"Color scale 'haxby' B. Haxby's color scheme for geoid & gravity [C=RGB -T-200/200]" \
    -I0.2 -By+l"mGal" -O -K >> $ps
# Add scale, directional rose
gmt psbasemap -R -J \
    -Lx14.0c/-2.3c+c50+w200k+l"Mercator projection. Scale (km)"+f \
    -UBL/-5p/-70p -O -K >> $ps
# Add coastlines, borders, rivers
gmt pscoast -R -J -P -Ia/thinnest,blue -Na -N1/thick,white -Wthin,darkslategray -Df -O -K >> $ps
# Texts
gmt pstext -R -J -N -O -K \
-F+jTL+f12p,21,white+jLB+a-315 >> $ps << EOF
286.7 5.8 Los Andes
EOF
# Add GMT logo
gmt logo -Dx7.0/-2.9+o0.1i/0.1i+w2c -O -K >> $ps
# Add subtitle
gmt pstext -R0/10/0/15 -JX10/10 -X0.0c -Y4.8c -N -O \
    -F+f10p,25,black+jLB >> $ps << EOF
1.0 13.6 Global gravity grid from CryoSat-2 and Jason-1, 1 min resolution, SIO, NOAA, NGA.
EOF
# Convert to image file using GhostScript
gmt psconvert Grav_VE_v.ps -A0.5c -E720 -Tj -Z
\end{lstlisting}

\subsection[\appendixname~\thesubsection]{GMT Script for Seismic Mapping}
\noindent\textbf{Listing A5:} GMT codes used to plot Figure~\ref{fig7} (seismic mapping).
\begin{lstlisting}[language=bash]
#!/bin/sh
# Purpose: seismicity map of Venezuela (earthquake distribution according to the IRIS database (1997-2021) of seismic events)
# GMT modules: gmtset, gmtdefaults, grdcut, makecpt, grdimage, psscale, grdcontour, psbasemap, gmtlogo, psconvert
# http://soliton.vm.bytemark.co.uk/pub/cpt-city/wkp/shadowxfox/tn/colombia.png.index.html
# GMT set up
gmt set FORMAT_GEO_MAP=dddF 
gmtdefaults -D > .gmtdefaults
# Extract a topographic subset of Venezuela and surroundings:
gmt grdcut ETOPO1_Ice_g_gmt4.grd -R286/300.5/0/12.5 -Gve_relief1.nc
gmt grdcut GEBCO_2019.nc -R286/300.5/0/12.5 -Gve_relief.nc
gdalinfo -stats ve_relief.nc
# Min=-4947.963 Max=5535.625
# Make color palette
gmt makecpt -Ccolombia.cpt > pauline.cpt
gmt makecpt -Cseis -T1.9/7.3/0.5 -Z > steps.cpt
ps=Seis_VE.ps
# Make raster image
gmt grdimage ve_relief.nc -Cpauline.cpt -R286/300.5/0/12.5 -JM6.5i -I+a15+ne0.75 -Xc -P -K > $ps
# Add legend
gmt psscale -Dg286/-1.0+w16.5c/0.4c+h+o0.0/0i+ml+e -R -J -Cpauline.cpt \
    -Bg500f50a500+l"Colormap..." \
    -I0.2 -By+lm -O -K >> $ps
# Add isolines
gmt grdcontour ve_relief1.nc -R -J -C500 -Wthinnest,brown -O -K >> $ps
# Add coastlines, borders, rivers
gmt pscoast -R -J -P \
    -Ia/thinner,blue -Na -N1/thickest,olivedrab1 -Wthin,lightcyan2 -Df -O -K >> $ps 
# Add grid
gmt psbasemap -R -J \
    --MAP_FRAME_AXES=WEsN \
    --FORMAT_GEO_MAP=ddd:mm:ssF \
    -Bpx4f2a2 -Bpyg4f2a2 -Bsxg4 -Bsyg2 \
    -B+t"Seismicity in Venezuela and Trinidad and Tobago: IRIS database (1997-2021)" -O -K >> $ps
# Add scale, directional rose
gmt psbasemap -R -J \
    -Lx14.0c/-3.6c+c50+w200k+l"Mercator projection. Scale: km"+f \
    -UBL/-5p/-110p -O -K >> $ps
# Add earthquake points
# separator in numbers of table: dot (.), not comma ! (British style)
gmt psxy -R -J quakes_VE.ngdc -Wfaint -i4,3,6,6s0.05 -h3 -Scc -Csteps.cpt -O -K >> $ps
# Add geological lines and points
gmt psxy -R -J volcanoes.gmt -St0.4c -Gred -Wthinnest -O -K >> $ps
# fabric and magnetic lineation picks fracture zones
gmt psxy -R -J GSFML_SF_FZ_KM.gmt -Wthicker,goldenrod1 -O -K >> $ps
gmt psxy -R -J GSFML_SF_FZ_RM.gmt -Wthicker,pink -O -K >> $ps
gmt psxy -R -J ridge.gmt -Sf0.5c/0.15c+l+t -Wthick,red -Gyellow -O -K >> $ps
gmt psxy -R -J ridge.gmt -Sc0.05c -Gred -Wthickest,red -O -K >> $ps
# tectonic plates
gmt psxy -R -J TP_Caribbean.txt -L -Wthickest,purple -O -K >> $ps
gmt psxy -R -J TP_Cocos.txt -L -Wthickest,purple -O -K >> $ps
gmt psxy -R -J TP_Nazca.txt -L -Wthickest,purple -O -K >> $ps
gmt psxy -R -J TP_South_Am.txt -L -Wthickest,purple -O -K >> $ps
# Texts
gmt pstext -R -J -N -O -K \
-F+f10p,17,black+jLB+a-0 >> $ps << EOF
293.20 10.58 Caracas
EOF
# Processed likewise for similar texts
gmt pslegend -R -J -Dx1.5/-3.0+w17.8c+o-2.0/0.1c \
    -F+pthin+ithinner+gwhite \
    --FONT=8p,black -O -K << FIN >> $ps
H 10 Helvetica Seismicity: earthquakes magnitude (M) from 1.9 to 7.3
N 9
S 0.3c c 0.3c red 0.01c 0.5c M (7.1-7.3)
S 0.3c c 0.3c tomato 0.01c 0.5c M (6.4-7.0)
S 0.3c c 0.3c orange 0.01c 0.5c M (5.6-6.3)
S 0.3c c 0.3c yellow 0.01c 0.5c M (4.8-5.5)
S 0.3c c 0.3c chartreuse1 0.01c 0.5c M (4.1-4.7)
S 0.3c c 0.3c chartreuse1 0.01c 0.5c M (3.4-4.0)
S 0.3c c 0.3c cyan3 0.01c 0.5c M (2.7-3.3)
S 0.3c c 0.3c blue 0.01c 0.5c M (1.9-2.6)
S 0.3c t 0.3c red 0.03c 0.5c Volcanoes
FIN
# Add GMT logo
gmt logo -Dx7.0/-4.5+o0.1i/0.1i+w2c -O -K >> $ps
# Add subtitle
gmt pstext -R0/10/0/15 -JX10/10 -X0.5c -Y5.9c -N -O \
    -F+f11p,25,black+jLB >> $ps << EOF
1.0 13.6 DEM: SRTM/GEBCO, 15 arc sec grid. Earthquakes: IRIS Seismic Event Database
EOF
# Convert to image file using GhostScript
gmt psconvert Seis_VE.ps -A1.7c -E720 -Tj -Z
\end{lstlisting}

%%%%%%%%%%%%%%%%%%%%%%%%%%%%%%%%%%%s%%%%%%%
\begin{adjustwidth}{-\extralength}{0cm}
%\printendnotes[custom] % Un-comment to print a list of endnotes

\reftitle{References}

% Please provide either the correct journal abbreviation (e.g. according to the “List of Title Word Abbreviations” http://www.issn.org/services/online-services/access-to-the-ltwa/) or the full name of the journal.
% Citations and References in Supplementary files are permitted provided that they also appear in the reference list here. 

%=====================================
% References, variant A: external bibliography
%=====================================
\begin{thebibliography}{999}

\bibitem[Laffaille \em{et~al.}(2009)Laffaille, Ferrer, and
  Laffaille]{Laffaille}
Laffaille, J.; Ferrer, C.; Laffaille, K.
\newblock {Venezuela: The Construction of Vulnerability and Its Relation to the
  High Seismic Risk}.
\newblock {\em Dev. Earth Surf. Process.} {\bf 2009}, {\em
  13},~99--114.
\newblock {{https://doi.org/10.1016/S0928-2025(08)10005-0}}.

\bibitem[Schmitz \em{et~al.}(2020)Schmitz, Hernández, Rocabado, Domínguez,
  Morales, Valleé, García, Sánchez-Rojas, Singer, Oropeza, Coronel, and
  Flores]{Schmitz2020}
Schmitz, M.; Hernández, J.J.; Rocabado, V.; Domínguez, J.; Morales, C.;
  Valleé, M.; García, K.; Sánchez-Rojas, J.; Singer, A.; Oropeza, J.;
  et~al.
\newblock {The Caracas, Venezuela, Seismic Microzoning Project: Methodology,
  results, and implementation for seismic risk reduction}.
\newblock {\em Prog. Disaster Sci.} {\bf 2020}, {\em 5},~100060.
\newblock {{https://doi.org/10.1016/j.pdisas.2019.100060}}.

\bibitem[Weber \em{et~al.}(2011)Weber, Saleh, Balkaransingh, Dixon, Ambeh,
  Leong, Rodriguez, and Miller]{Weber}
Weber, J.C.; Saleh, J.; Balkaransingh, S.; Dixon, T.; Ambeh, W.; Leong, T.;
  Rodriguez, A.; Miller, K.
\newblock {Triangulation-to-GPS and GPS-to-GPS geodesy in Trinidad, West
  Indies: Neotectonics, seismic risk, and geologic implications}.
\newblock {\em Mar. Pet. Geol.} {\bf 2011}, {\em 28},~200--211.
\newblock {{https://doi.org/10.1016/j.marpetgeo.2009.07.010}}.

\bibitem[Khan \em{et~al.}(2020)Khan, Gupta, and Gupta]{KHAN2020101642}
Khan, A.; Gupta, S.; Gupta, S.K.
\newblock Multi-hazard disaster studies: Monitoring, detection, recovery, and
  management, based on emerging technologies and optimal techniques.
\newblock {\em Int. J. Disaster Risk Reduct.} {\bf 2020},
  {\em 47},~101642.
\newblock {{https://doi.org/10.1016/j.ijdrr.2020.101642}}.

\bibitem[Mesgar and Jalilvand(2017)]{MESGAR2017454}
Mesgar, M.A.A.; Jalilvand, P.
\newblock Vulnerability Analysis of the Urban Environments to Different Seismic
  Scenarios: Residential Buildings and Associated Population Distribution
  Modelling through Integrating Dasymetric Mapping Method and GIS.
\newblock {\em Procedia Eng.} {\bf 2017}, {\em 198},~454--466.
\newblock 
  {{https://doi.org/10.1016/j.proeng.2017.07.100}}.

\bibitem[Nath(2005)]{NATH2005329}
Nath, S.K.
\newblock An initial model of seismic microzonation of Sikkim Himalaya through
  thematic mapping and GIS integration of geological and strong motion
  features.
\newblock {\em J. Asian Earth Sci.} {\bf 2005}, {\em 25},~329--343.
\newblock {{https://doi.org/10.1016/j.jseaes.2004.03.002}}.

\bibitem[Kwag \em{et~al.}(2022)Kwag, Choi, Hahm, Eem, and Ju]{KWAG20227230}
Kwag, S.; Choi, E.; Hahm, D.; Eem, S.; Ju, B.S.
\newblock On improving seismic risk and cost for nuclear energy facility based
  on multi-objective optimization considering seismic correlation.
\newblock {\em Energy Rep.} {\bf 2022}, {\em 8},~7230--7241.
\newblock {{https://doi.org/10.1016/j.egyr.2022.05.241}}.

\bibitem[Kwag \em{et~al.}(2021)Kwag, Choi, Eem, Ha, and Hahm]{KWAG2021107809}
Kwag, S.; Choi, E.; Eem, S.; Ha, J.G.; Hahm, D.
\newblock Toward improvement of sampling-based seismic probabilistic safety
  assessment method for nuclear facilities using composite distribution and
  adaptive discretization.
\newblock {\em Reliab. Eng. Syst. Saf.} {\bf 2021}, {\em
  215},~107809.
\newblock {{https://doi.org/10.1016/j.ress.2021.107809}}.

\bibitem[Liu \em{et~al.}(2011)Liu, Cui, Yuan, Wang, Jin, and Wang]{LIU20111084}
Liu, H.; Cui, X.; Yuan, D.; Wang, Z.; Jin, J.; Wang, M.
\newblock Study of Earthquake Disaster Population Risk Based on GIS A Case
  Study of Wenchuan Earthquake Region.
\newblock {\em Procedia Environ. Sci.} {\bf 2011}, {\em
  11},~1084--1091.
\newblock 2011 2nd International Conference on Challenges in Environmental
  Science and Computer Engineering (CESCE 2011),
  {{https://doi.org/10.1016/j.proenv.2011.12.164}}.

\bibitem[Zheng \em{et~al.}(2022)Zheng, Feng, and Ishiwatari]{ijerph191912351}
Zheng, X.; Feng, C.; Ishiwatari, M.
\newblock Examining the Indirect Death Surveillance System of The Great East
  Japan Earthquake and Tsunami.
\newblock {\em Int. J. Environ. Res. Public Health} {\bf 2022}, {\em 19}, 12351.
\newblock {{https://doi.org/10.3390/ijerph191912351}}.

\bibitem[Peng \em{et~al.}(2022)Peng, Jiang, Ma, Su, Cai, and Zheng]{rs14174269}
Peng, C.; Jiang, P.; Ma, Q.; Su, J.; Cai, Y.; Zheng, Y.
\newblock Chinese Nationwide Earthquake Early Warning System and Its
  Performance in the 2022 Lushan M6.1 Earthquake.
\newblock {\em Remote Sens.} {\bf 2022}, {\em 14}, 4269.
\newblock {{https://doi.org/10.3390/rs14174269}}.

\bibitem[Adu-Gyamfi and Shaw(2022)]{ijerph191811469}
Adu-Gyamfi, B.; Shaw, R.
\newblock Risk Awareness and Impediments to Disaster Preparedness of Foreign
  Residents in the Tokyo Metropolitan Area, Japan.
\newblock {\em Int. J. Environ. Res. Public Health} {\bf 2022}, {\em 19}, 11469.
\newblock {{https://doi.org/10.3390/ijerph191811469}}.

\bibitem[Yousuf \em{et~al.}(2020)Yousuf, Bukhari, Bhat, and
  Ali]{YOUSUF2020100064}
Yousuf, M.; Bukhari, S.K.; Bhat, G.R.; Ali, A.
\newblock Understanding and managing earthquake hazard visa viz disaster
  mitigation strategies in Kashmir valley, NW Himalaya.
\newblock {\em Prog. Disaster Sci.} {\bf 2020}, {\em 5},~100064.
\newblock {{https://doi.org/10.1016/j.pdisas.2020.100064}}.

\bibitem[Kamata \em{et~al.}(2022)Kamata, Seto, Suppasri, Sasaki, Egawa, and
  Imamura]{KAMATA2022103253}
Kamata, H.; Seto, S.; Suppasri, A.; Sasaki, H.; Egawa, S.; Imamura, F.
\newblock A study on hypothermia and associated countermeasures in tsunami
  disasters: A case study of Miyagi Prefecture during the 2011 great East Japan
  earthquake.
\newblock {\em Int. J. Disaster Risk Reduct.} {\bf 2022},
  {\em 81},~103253.
\newblock {{https://doi.org/10.1016/j.ijdrr.2022.103253}}.

\bibitem[Shao \em{et~al.}(2022)Shao, Wu, Ji, Diao, Shi, Li, Long, Zhang, Wang,
  Wei, Wang, Liu, Liu, Pan, Yin, Liu, Feng, Zou, Cheng, Lu, Xu, and
  Li]{SHAO2022100177}
Shao, Z.; Wu, Y.; Ji, L.; Diao, F.; Shi, F.; Li, Y.; Long, F.; Zhang, H.; Wang,
  W.; Wei, W.;  et~al.
\newblock Assessment of strong earthquake risk in the Chinese mainland from
  2021 to 2030.
\newblock {\em Earthq. Res. Adv.} {\bf 2022}, {\em 2},~100177. %MDPI: please add the volume. <-- Answer: added (vol. 2).
\newblock {{https://doi.org/10.1016/j.eqrea.2022.100177}}.

\bibitem[Yaghmaei-Sabegh \em{et~al.}(2022)Yaghmaei-Sabegh, Pavel, Shahvar, and
  Qadri]{YAGHMAEISABEGH2022107173}
Yaghmaei-Sabegh, S.; Pavel, F.; Shahvar, M.; Qadri, S.T.
\newblock Empirical frequency content models based on intermediate-depth
  earthquake ground-motions.
\newblock {\em Soil Dyn. Earthq. Eng.} {\bf 2022}, {\em
  155},~107173.
\newblock {{https://doi.org/10.1016/j.soildyn.2022.107173}}.

\bibitem[Wessel \em{et~al.}(2019)Wessel, Luis, Uieda, Scharroo, Wobbe, Smith,
  and Tian]{Wessel}
Wessel, P.; Luis, J.F.; Uieda, L.; Scharroo, R.; Wobbe, F.; Smith, W.H.F.;
  Tian, D.
\newblock {The Generic Mapping Tools version 6}.
\newblock {\em Geochem. Geophys. Geosystems} {\bf 2019}, {\em
  20},~5556--5564.
\newblock {{https://doi.org/10.1029/2019GC008515}}.

\bibitem[Castillo \em{et~al.}(2011)Castillo, López-Almansa, and
  Pujades]{Castillo}
Castillo, A.; López-Almansa, F.; Pujades, L.G.
\newblock {Seismic risk analysis of urban non-engineered buildings: application
  to an informal settlement in Mérida, Venezuela}.
\newblock {\em Nat. Hazards} {\bf 2011}, {\em 59},~891--916.
\newblock {{https://doi.org/10.1007/s11069-011-9805-9}}.

\bibitem[Salcedo(2009)]{Salcedo}
Salcedo, D.A.
\newblock {Behavior of a landslide prior to inducing a viaduct failure,
  Caracas–La Guaira highway, Venezuela}.
\newblock {\em Eng. Geol.} {\bf 2009}, {\em 109},~16--30.
\newblock {{https://doi.org/10.1016/j.enggeo.2009.02.001}}.

\bibitem[Rodríguez \em{et~al.}(2018)Rodríguez, Diederix, Torres, Audemard,
  Hernández, Singer, Bohórquez, and Yepez]{Rodriguez2018}
Rodríguez, L.; Diederix, H.; Torres, E.; Audemard, F.; Hernández, C.; Singer,
  A.; Bohórquez, O.; Yepez, S.
\newblock {Identification of the seismogenic source of the 1875 Cucuta
  earthquake on the basis of a combination of neotectonic, paleoseismologic and
  historic seismicity studies}.
\newblock {\em J. S. Am. Earth Sci.} {\bf 2018}, {\em
  82},~274--291.
\newblock {{https://doi.org/10.1016/j.jsames.2017.09.019}}.

\bibitem[{Pousse Beltran} \em{et~al.}(2016){Pousse Beltran}, Pathier, Jouanne,
  Vassallo, Reinoza, Audemard, Doin, and Volat]{Pousse2016}
{Pousse Beltran}, L.; Pathier, E.; Jouanne, F.; Vassallo, R.; Reinoza, C.;
  Audemard, F.; Doin, M.P.; Volat, M.
\newblock {Spatial and temporal variations in creep rate along the El Pilar
  fault at the Caribbean-South American plate boundary (Venezuela), from
  InSAR}.
\newblock {\em J. Geophys. Res. Solid Earth} {\bf 2016}, {\em
  121},~8276--8296.
\newblock {{https://doi.org/10.1002/2016JB013121}}.

\bibitem[Lyu \em{et~al.}(2018)Lyu, Shen, and Arulrajah]{su10020304}
Lyu, H.M.; Shen, J.S.; Arulrajah, A.
\newblock Assessment of Geohazards and Preventative Countermeasures Using AHP
  Incorporated with GIS in Lanzhou, China.
\newblock {\em Sustainability} {\bf 2018}, {\em 10}.
\newblock {{https://doi.org/10.3390/su10020304}}.

\bibitem[Karymbalis \em{et~al.}(2021)Karymbalis, Andreou, Batzakis, Tsanakas,
  and Karalis]{su131810232}
Karymbalis, E.; Andreou, M.; Batzakis, D.V.; Tsanakas, K.; Karalis, S.
\newblock Integration of GIS-Based Multicriteria Decision Analysis and Analytic
  Hierarchy Process for Flood-Hazard Assessment in the Megalo Rema River
  Catchment (East Attica, Greece).
\newblock {\em Sustainability} {\bf 2021}, {\em 13}.
\newblock {{https://doi.org/10.3390/su131810232}}.

\bibitem[Lemenkova(2022)]{Lemenkova202212}
Lemenkova, P.
\newblock {Handling Dataset with Geophysical and Geological Variables on the
  Bolivian Andes by the GMT Scripts}.
\newblock {\em Data} {\bf 2022}, {\em 7},~74.
\newblock {{https://doi.org/10.3390/data7060074}}.

\bibitem[He \em{et~al.}(2021)He, Bai, Shi, Li, Qi, and Li]{app11219919}
He, B.; Bai, M.; Shi, H.; Li, X.; Qi, Y.; Li, Y.
\newblock Risk Assessment of Pipeline Engineering Geological Disaster Based on
  GIS and WOE-GA-BP Models.
\newblock {\em Appl. Sci.} {\bf 2021}, {\em 11}, 9919.
\newblock {{https://doi.org/10.3390/app11219919}}.

\bibitem[Valkanou \em{et~al.}(2021)Valkanou, Karymbalis, Papanastassiou,
  Soldati, Chalkias, and Gaki-Papanastassiou]{ijgi10030118}
Valkanou, K.; Karymbalis, E.; Papanastassiou, D.; Soldati, M.; Chalkias, C.;
  Gaki-Papanastassiou, K.
\newblock Assessment of Neotectonic Landscape Deformation in Evia Island,
  Greece, Using GIS-Based Multi-Criteria Analysis.
\newblock {\em ISPRS Int. J. -Geo-Inf.} {\bf 2021}, {\em
  10}, 118.
\newblock {{https://doi.org/10.3390/ijgi10030118}}.

\bibitem[González-Longatt \em{et~al.}(2015)González-Longatt, Medina, and
  González]{Gonzalez}
González-Longatt, F.; Medina, H.; González, J.S.
\newblock {Spatial interpolation and orographic correction to estimate wind
  energy resource in Venezuela}.
\newblock {\em Renew. Sustain. Energy Rev.} {\bf 2015}, {\em
  48},~1--16.
\newblock {{https://doi.org/10.1016/j.rser.2015.03.042}}.

\bibitem[Hackley \em{et~al.}(2006)Hackley, Urbani, Karlsen, and
  Garrity]{Hackley}
Hackley, P.; Urbani, F.; Karlsen, A.W.; Garrity, C.P.
\newblock Geologic shaded relief map of Venezuela.
\newblock U.S. Geological Survey, Open File Report 2005–1038,  2006.
\newblock  Available online: \url{https://pubs.usgs.gov/of/2005/1038/index.htm} (accessed on 1 November 2022). %MDPI: Please add the access date (Format: Date Month Year). e.g., (accessed on 1 January 2020). <-- Answer: added information on access date (1 November 2022).


\bibitem[Schmitz \em{et~al.}(2021)Schmitz, Ramírez, Mazuera, Ávila, Yegres,
  Bezada, and Levander]{Schmitz2021}
Schmitz, M.; Ramírez, K.; Mazuera, F.; Ávila, J.; Yegres, L.; Bezada, M.;
  Levander, A.
\newblock {Moho depth map of northern Venezuela based on wide-angle seismic
  studies}.
\newblock {\em J. South Am. Earth Sci.} {\bf 2021}, {\em
  107},~103088.
\newblock {{https://doi.org/10.1016/j.jsames.2020.103088}}.

\bibitem[Sun \em{et~al.}(2008)Sun, Chun, Ha, Chung, and Kim]{Sun}
Sun, C.G.; Chun, S.H.; Ha, T.G.; Chung, C.K.; Kim, D.S.
\newblock {Development and application of a GIS-based tool for
  earthquake-induced hazard prediction}.
\newblock {\em Comput. Geotech.} {\bf 2008}, {\em 35},~436--449.
\newblock {{https://doi.org/10.1016/j.compgeo.2007.08.001}}.

\bibitem[Escalona and Mann(2011)]{ESCALONA20118}
Escalona, A.; Mann, P.
\newblock Tectonics, basin subsidence mechanisms, and paleogeography of the
  Caribbean-South American plate boundary zone. Thematic Set on: Tectonics, basinal framework, and petroleum systems
  of eastern Venezuela, the Leeward Antilles, Trinidad and Tobago, and offshore
  areas.
\newblock {\em Mar. Pet. Geol.} {\bf 2011}, {\em 28},~8--39.
\newblock  {{https://doi.org/10.1016/j.marpetgeo.2010.01.016}}.

\bibitem[Dhont \em{et~al.}(2012)Dhont, Monod, Hervouët, Backé, Klarica, and
  Choy]{DHONT2012245}
Dhont, D.; Monod, B.; Hervouët, Y.; Backé, G.; Klarica, S.; Choy, J.E.
\newblock 3D geological modeling of the Trujillo block: Insights for crustal
  escape models of the Venezuelan Andes.
\newblock {\em J. Am. Earth Sci.} {\bf 2012}, {\em
  39},~245--251.
\newblock {https://doi.org/10.1016/j.jsames.2012.04.003}.

\bibitem[Masy \em{et~al.}(2011)Masy, Niu, Levander, and Schmitz]{Masy}
Masy, J.; Niu, F.; Levander, A.; Schmitz, M.
\newblock {Mantle flow beneath northwestern Venezuela: Seismic evidence for a
  deep origin of the Mérida Andes}.
\newblock {\em Earth Planet. Sci. Lett.} {\bf 2011}, {\em
  305},~396--404.
\newblock {{https://doi.org/10.1016/j.epsl.2011.03.024}}.

\bibitem[Pérez \em{et~al.}(2018)Pérez, Wesnousky, {De La Rosa}, Márquez,
  Uzcátegui, Quintero, Liberal, {Mora-Páez}, and Szeliga]{Perez2018}
Pérez, O.J.; Wesnousky, S.G.; {De La Rosa}, R.; Márquez, J.; Uzcátegui, R.;
  Quintero, C.; Liberal, L.; {Mora-Páez}, H.; Szeliga, W.
\newblock {On the interaction of the North Andes plate with the Caribbean and
  South American plates in northwestern South America from GPS geodesy and
  seismic data}.
\newblock {\em Geophys. J. Int.} {\bf 2018}, {\em
  214},~1986--2001.
\newblock {{https://doi.org/10.1093/gji/ggy230}}.

\bibitem[Suárez and Nábělek(1990)]{Suarez}
Suárez, G.; Nábělek, J.
\newblock {The 1967 Caracas Earthquake: Fault geometry, direction of rupture
  propagation and seismotectonic implications}.
\newblock {\em J. Geophys. Res.} {\bf 1990}, {\em
  95},~17459--17474.
\newblock {{https://doi.org/10.1029/JB095iB11p17459}}.

\bibitem[Audemard(1993)]{Audemard1993}
Audemard, F.A.
\newblock {Néotectonique, Sismotectonique et aléa Sismique du nord-ouest
  du Vénézuela (Système de Failles d'Oca-Ancón)}; PhD Thesis, University of
  Montpellier I, Montpellier, France,  1993;
\newblock 369p.

\bibitem[Backé \em{et~al.}(2006)Backé, Dhont, and Hervouët]{Backe}
Backé, G.; Dhont, D.; Hervouët, Y.
\newblock {Spatial and temporal relationships between compression, strike-slip
  and extension in the Central Venezuelan Andes: Clues for Plio-Quaternary
  tectonic escape}.
\newblock {\em Tectonophysics} {\bf 2006}, {\em 425},~25--53.
\newblock {{https://doi.org/10.1016/j.tecto.2006.06.005}}.

\bibitem[Marshall and Russo(2005)]{Marshall}
Marshall, J.L.; Russo, R.M.
\newblock {Relocated aftershocks of the March 10, 1988 Trinidad earthquake:
  Normal faulting, slab detachment and extension at upper mantle depths}.
\newblock {\em Tectonophysics} {\bf 2005}, {\em 398},~101--114.
\newblock {{https://doi.org/10.1016/j.tecto.2004.11.010}}.

\bibitem[Reinoza \em{et~al.}(2015)Reinoza, Jouanne, Audemard, Schmitz, and
  Beck]{Reinoza}
Reinoza, C.; Jouanne, F.; Audemard, F.A.; Schmitz, M.; Beck, C.
\newblock {Geodetic exploration of strain along the El Pilar Fault in
  northeastern Venezuela}.
\newblock {\em J. Geophys. Res. Solid Earth} {\bf 2015}, {\em
  120},~1993--2013.
\newblock {{https://doi.org/10.1002/2014JB011483}}.

\bibitem[Harbitz \em{et~al.}(2012)Harbitz, Glimsdal, Bazin, Zamora, Løvholt,
  Bungum, Smebye, Gauer, and Kjekstad]{Harbitz}
Harbitz, C.B.; Glimsdal, S.; Bazin, S.; Zamora, N.; Løvholt, F.; Bungum, H.;
  Smebye, H.; Gauer, P.; Kjekstad, O.
\newblock {Tsunami hazard in the Caribbean: Regional exposure derived from
  credible worst case scenarios}.
\newblock {\em Cont. Shelf Res.} {\bf 2012}, {\em 38},~1--23.
\newblock {{https://doi.org/10.1016/j.csr.2012.02.006}}.

\bibitem[Audemard(2007)]{Audemard}
Audemard, F.A.
\newblock {Revised seismic history of the El Pilar fault, Northeastern
  Venezuela, from the Cariaco 1997 earthquake and recent preliminary
  paleoseismic results}.
\newblock {\em J. Seismol.} {\bf 2007}, {\em 11},~311--326.
\newblock {{https://doi.org/10.1007/s10950-007-9054-2}}.

\bibitem[Pérez \em{et~al.}(1997)Pérez, Sanz, and Lagos]{Perez}
Pérez, O.J.; Sanz, C.; Lagos, G.
\newblock {Microseismicity, tectonics and seismic potential in southern
  Caribbean and northern Venezuela}.
\newblock {\em J. Seismol.} {\bf 1997}, {\em 1},~15--28.
\newblock {{https://doi.org/10.1023/A:1009710122083}}.

\bibitem[Barr and Robson(1963)]{Barr}
Barr, K.G.; Robson, G.R.
\newblock {Seismic Delays in the Eastern Caribbean}.
\newblock {\em Geophys. J. Int.} {\bf 1963}, {\em
  7},~342--349.
\newblock {{https://doi.org/10.1111/j.1365-246X.1963.tb05555.x}}.

\bibitem[Hayes \em{et~al.}(2014)Hayes, Mcnamara, Seidman, and Roger]{Hayes}
Hayes, G.P.; Mcnamara, D.E.; Seidman, L.; Roger, J.
\newblock {Quantifying potential earthquake and tsunami hazard in the Lesser
  Antilles subduction zone of the Caribbean region}.
\newblock {\em Geophys. J. Int.} {\bf 2014}, {\em
  196},~510--521.
\newblock {{https://doi.org/10.1093/gji/ggt385}}.

\bibitem[Audemard(2006)]{Audemard2006}
Audemard, F.A.
\newblock {Surface rupture of the Cariaco July 09, 1997 earthquake on the El
  Pilar fault, northeastern Venezuela}.
\newblock {\em Tectonophysics} {\bf 2006}, {\em 424},~19--39.
\newblock {{https://doi.org/10.1016/j.tecto.2006.04.018}}.

\bibitem[Baumbach \em{et~al.}(2004)Baumbach, Grosser, {Romero Torres}, {Rojas
  Gonzales}, Sobiesiak, and Welle]{Baumbach}
Baumbach, M.; Grosser, H.; {Romero Torres}, G.; {Rojas Gonzales}, J.L.;
  Sobiesiak, M.; Welle, W.
\newblock {Aftershock pattern of the July 9, 1997 Mw=6.9 Cariaco earthquake in
  Northeastern Venezuela}.
\newblock {\em Tectonophysics} {\bf 2004}, {\em 379},~1--23.
\newblock {{https://doi.org/10.1016/j.tecto.2003.10.018}}.

\bibitem[Colón \em{et~al.}(2015)Colón, Audemard, Beck, Avila, Padrón, {De
  Batist}, Paolini, Leal, and {Van Welden}]{Colon}
Colón, S.; Audemard, F.A.; Beck, C.; Avila, J.; Padrón, C.; {De Batist}, M.;
  Paolini, M.; Leal, A.F.; {Van Welden}, A.
\newblock {The 1900 Mw 7.6 earthquake offshore north–central Venezuela: Is La
  Tortuga or San Sebastián the source fault?}
\newblock {\em Mar. Pet. Geol.} {\bf 2015}, {\em 67},~498--511.
\newblock {{https://doi.org/10.1016/j.marpetgeo.2015.06.005}}.

\bibitem[Pousse-Beltran \em{et~al.}(2018)Pousse-Beltran, Vassallo, Audemard,
  Jouanne, Oropeza, Garambois, and Aray]{Pousse2018}
Pousse-Beltran, L.; Vassallo, R.; Audemard, F.; Jouanne, F.; Oropeza, J.;
  Garambois, S.; Aray, J.
\newblock {Earthquake geology of the last millennium along the Boconó Fault,
  Venezuela}.
\newblock {\em Tectonophysics} {\bf 2018}, {\em 747--748},~40--53.
\newblock {{https://doi.org/10.1016/j.tecto.2018.09.010}}.

\bibitem[Rial(1978)]{Rial}
Rial, J.A.
\newblock {The Caracas, Venezuela, earthquake of July 1967: A multiple-source
  event}.
\newblock {\em J. Geophys. Res. Solid Earth} {\bf 1978}, {\em
  83},~5405--5414.
\newblock {{https://doi.org/10.1029/JB083iB11p05405}}.

\bibitem[Yamazaki \em{et~al.}(2005)Yamazaki, Audemard, Altez, Hernández,
  Orihuela, Safina, Schmitz, Tanaka, and Kagawa]{Yamazaki}
Yamazaki, Y.; Audemard, F.; Altez, R.; Hernández, J.; Orihuela, N.; Safina,
  S.; Schmitz, M.; Tanaka, I.; Kagawa, H.
\newblock {Estimation of the seismic intensity in Caracas during the 1812
  earthquake using seismic microzonation methodology}.
\newblock {\em Rev. GeográFica Venez.} {\bf 2005}, {\em Número
  Especial},~199--216.

\bibitem[Deville and Guerlais(2009)]{Deville}
Deville, E.; Guerlais, S.H.
\newblock {Cyclic activity of mud volcanoes: Evidences from Trinidad (SE
  Caribbean)}.
\newblock {\em Mar. Pet. Geol.} {\bf 2009}, {\em 26},~1681--1691.
\newblock {{https://doi.org/10.1016/j.marpetgeo.2009.03.002}}.

\bibitem[Latchman \em{et~al.}(1996)Latchman, Ambeh, and Lynch]{Latchman}
Latchman, J.L.; Ambeh, W.B.; Lynch, L.L.
\newblock {Attenuation of seismic waves in the Trinidad and Tobago area}.
\newblock {\em Tectonophysics} {\bf 1996}, {\em 253},~111--127.
\newblock {{https://doi.org/10.1016/0040-1951(95)00050-X}}.

\bibitem[Bucci and Tuttle(2022)]{Bucci}
Bucci, M.G.; Tuttle, M.P.
\newblock {Liquefaction Susceptibility and Hazard Mapping: Insights From Case
  Histories of Earthquake-Induced Liquefaction}.
\newblock {\em Ref. Modul. Earth Syst. Environ. Sci.}
  {\bf 2022}, {\em 2},~563--582.
\newblock {{https://doi.org/10.1016/B978-0-12-818234-5.00054-7}}.

\bibitem[Diebold \em{et~al.}(1981)Diebold, Stoffa, Buhl, and Truchan]{Diebold}
Diebold, J.B.; Stoffa, P.L.; Buhl, P.; Truchan, M.
\newblock {Venezuela Basin crustal structure}.
\newblock {\em J. Geophys. Res. Solid Earth} {\bf 1981}, {\em
  86},~7901–7923.
\newblock {{https://doi.org/10.1029/JB086iB09p07901}}.

\bibitem[Fajardo \em{et~al.}(2020)Fajardo, Aubourg, Nivière, Regard, and
  Uzcátegui]{Fajardo}
Fajardo, A.; Aubourg, C.; Nivière, B.; Regard, V.; Uzcátegui, R.
\newblock {Neotectonic evolution of Northeastern Venezuela. Geomorphological
  and seismic analysis in the Monagas Fold and Thrust Belt}.
\newblock {\em J. S. Am. Earth Sci.} {\bf 2020}, {\em
  104},~102867.
\newblock {{https://doi.org/10.1016/j.jsames.2020.102867}}.

\bibitem[Lemenkova(2020)]{Lemenkova2020b}
Lemenkova, P.
\newblock {GRASS GIS for topographic and geophysical mapping of the Peru-Chile
  Trench}.
\newblock {\em Forum Geogr.} {\bf 2020}, {\em 19},~143--157.
\newblock {{https://doi.org/10.5775/fg.2020.009.d}}.

\bibitem[Pérez and Aggarwal(1981)]{Perez1981}
Pérez, O.J.; Aggarwal, Y.P.
\newblock {Present-day tectonics of the southeastern Caribbean and northeastern
  Venezuela}.
\newblock {\em J. Geophys. Res. Solid Earth} {\bf 1981}, {\em
  86},~10791--10804.
\newblock {{https://doi.org/10.1029/JB086iB11p10791}}.

\bibitem[Avila-García \em{et~al.}(2022)Avila-García, Schmitz,
  Mortera-Gutierrez, Bandy, Yegres, Zelt, and
  Aray-Castellano]{AVILAGARCIA2022103853}
Avila-García, J.; Schmitz, M.; Mortera-Gutierrez, C.; Bandy, W.; Yegres, L.;
  Zelt, C.; Aray-Castellano, J.
\newblock Crustal structure and tectonic implications of the southernmost
  Merida Andes, Venezuela, from wide-angle seismic data analysis.
\newblock {\em J. South Am. Earth Sci.} {\bf 2022}, {\em
  116},~103853.
\newblock {{https://doi.org/10.1016/j.jsames.2022.103853}}.

\bibitem[Gough and Heirtzler(1969)]{Gough}
Gough, D.I.; Heirtzler, J.R.
\newblock {Magnetic Anomalies and Tectonics of the Cayman Trough}.
\newblock {\em Geophys. J. Int.} {\bf 1969}, {\em 18},~33--49.
\newblock {{https://doi.org/10.1111/j.1365-246X.1969.tb00260.x}}.

\bibitem[Lemenkova(2020)]{Lemenkova2020a}
Lemenkova, P.
\newblock {Geomorphology of the Puerto Rico Trench and Cayman Trough in the
  Context of the Geological Evolution of the Caribbean Sea}.
\newblock {\em Ann. Univ. Mariae-Curie-Sklodowska, Sect. Geogr. Geol. Mineral. Petrogr.} {\bf 2020}, {\em
  75},~115--141.
\newblock {{https://doi.org/10.17951/b.2020.75.115-141}}.

\bibitem[Bezada \em{et~al.}(2008)Bezada, Schmitz, Jácome, Rodríguez,
  Audemard, and Izarra]{Bezada}
Bezada, M.J.; Schmitz, M.; Jácome, M.I.; Rodríguez, J.; Audemard, F.; Izarra,
  C.
\newblock {Crustal structure in the Falcón Basin area, northwestern Venezuela,
  from seismic and gravimetric evidence}.
\newblock {\em J. Geodyn.} {\bf 2008}, {\em 45},~191--200.
\newblock {{https://doi.org/10.1016/j.jog.2007.11.002}}.

\bibitem[Mazuera \em{et~al.}(2019)Mazuera, Schmitz, Escalona, Zelt, and
  Levander]{Mazuera}
Mazuera, F.; Schmitz, M.; Escalona, A.; Zelt, C.; Levander, A.
\newblock {Lithospheric structure of northwestern Venezuela from wide-angle
  seismic data: Implications for the understanding of continental margin
  evolution}.
\newblock {\em J. Geophys. Res. Solid Earth} {\bf 2019}, {\em
  124},~13124--13149.
\newblock {{https://doi.org/10.1029/2019JB017892}}.

\bibitem[Group(2020)]{GEBCO}
Group, G.C.
\newblock GEBCO 2020 Grid, 2020.
\newblock  Available online: %MDPI: please provide a new website link, doi can not be used as a link. <-- Answer: added URL (https://www.bodc.ac.uk/data/published_data_library/catalogue/10.5285/a29c5465-b138-234d-e053-6c86abc040b9/)
 \url{https://www.bodc.ac.uk/data/published_data_library/catalogue/10.5285/a29c5465-b138-234d-e053-6c86abc040b9/} (accessed on 1 November 2022). %MDPI: Please add the access date (Format: Date Month Year). e.g., (accessed on 1 January 2020). <-- Answer: added access date (1 November 2022).


\bibitem[Schenke(2016)]{Schenke}
Schenke, H., {General Bathymetric Chart of the Oceans (GEBCO)}.
\newblock In {\em {Encyclopedia of Earth Sciences Series}}; Chapter Encyclopedia of Marine Geosciences; Springer:
  Dordrecht, The Netherlands,  2016; 
  pp. 268--269.
\newblock {{https://doi.org/10.1007/978-94-007-6238-1\_63}}.

\bibitem[Mcmichael-Phillips(2021)]{Mcmichael}
Mcmichael-Phillips, J.
\newblock {The Nippon Foundation-GEBCO Seabed 2030 Project: The Most Ambitious
  Seafloor Mapping Initiative in History}.
\newblock {\em Mar. Technol. Soc. J.} {\bf 2021}, {\em 55},~25--28.
\newblock {{https://doi.org/10.4031/MTSJ.55.3.3}}.

\bibitem[Cortina \em{et~al.}(2019)Cortina, Wigley, and {Pe'eri}]{Cortina}
Cortina, G.C.Z.; Wigley, R.; {Pe'eri}, S. {The GEBCO and NOAA Chart Adequacy
  Workshop}.
\newblock In {\em {Abstracts of the International Cartographic Conference
  (ICA)}}; \hl{ICA:} %MDPI: Please add the location of the publisher (City and Country).
  2019; Chapter~1, p.~51.
\newblock https://doi.org/10.5194/ica-abs-1-51-2019.

\bibitem[Lemenkova(2022)]{Lemenkova2022a}
Lemenkova, P.
\newblock {Tanzania Craton, Serengeti Plain and Eastern Rift Valley: mapping of
  geospatial data by scripting techniques}.
\newblock {\em Est. J. Earth Sci.} {\bf 2022}, {\em
  71},~61--79.
\newblock {{https://doi.org/10.3176/earth.2022.05}}.

\bibitem[Pavlis \em{et~al.}(2012)Pavlis, Holmes, Kenyon, and Factor]{Pavlis}
Pavlis, N.K.; Holmes, S.; Kenyon, S.C.; Factor, J.K.
\newblock {The development and evaluation of the Earth Gravitational Model 2008
  (EGM2008)}.
\newblock {\em J. Geophys. Res.} {\bf 2012}, {\em 117},~B04406.
\newblock {{https://doi.org/10.1029/2011JB008916}}.

\bibitem[Sandwell \em{et~al.}(2014)Sandwell, Müller, Smith, Garcia, and
  Francis]{Sandwell}
Sandwell, D.T.; Müller, R.D.; Smith, W.H.F.; Garcia, E.; Francis, R.
\newblock {New global marine gravity model from CryoSat-2 and Jason-1 reveals
  buried tectonic structure}.
\newblock {\em Science} {\bf 2014}, {\em 7346},~65--67.
\newblock {{https://doi.org/10.1126/science.1258213}}.

\bibitem[Lemenkova(2022{\natexlab{a}})]{Lemenkova2022b}
Lemenkova, P.
\newblock {Console-Based Mapping of Mongolia Using GMT Cartographic Scripting
  Toolset for Processing TerraClimate Data}.
\newblock {\em Geosciences} {\bf 2022}, {\em 12},~140.
\newblock {{https://doi.org/10.3390/geosciences12030140}}.

\bibitem[Lemenkova(2022{\natexlab{b}})]{Lemenkova202217}
Lemenkova, P.
\newblock {Mapping Climate Parameters over the Territory of Botswana Using GMT
  and Gridded Surface Data from TerraClimate}.
\newblock {\em ISPRS Int. J. -Geo-Inf.} {\bf 2022}, {\em
  11},~473.
\newblock {{https://doi.org/10.3390/ijgi11090473}}.

\bibitem[Lemenkova(2021{\natexlab{a}})]{Lemenkova2021b}
Lemenkova, P.
\newblock {Topography of the Aleutian Trench south-east off Bowers Ridge,
  Bering Sea, in the context of the geological development of North Pacific
  Ocean}.
\newblock {\em Baltica} {\bf 2021}, {\em 34},~27--46.
\newblock {{https://doi.org/10.5200/baltica.2021.1.3}}.

\bibitem[Lemenkova(2021{\natexlab{b}})]{Lemenkova2021d}
Lemenkova, P.
\newblock {The visualization of geophysical and geomorphologic data from the
  area of Weddell Sea by the Generic Mapping Tools}.
\newblock {\em Stud. Quat.} {\bf 2021}, {\em 38},~19--32.
\newblock {{https://doi.org/10.24425/sq.2020.133759}}.

\bibitem[Callejón and {von der Dick}(2002)]{Callejon}
Callejón, A.F.; {von der Dick}, H.H. {Numerical identification of microseeps
  insurface soil gases of western Venezuela, and its significance for
  hydrocarbon exploration}.
\newblock In {\em {Surface Exploration Case Histories: Applications of
  Geochemistry, Magnetics, and Remote Sensing}}; AAPG Studies in Geology 48 and SEG Geophysical References Series; AAPG: \hl{Tulsa, OK, USA,} 2002; Volume~11, \mbox{pp.
  381--392.}

\bibitem[De~Rodriguez(1951)]{Rodriguez}
De~Rodriguez, G.
\newblock {Fundamental Geological Characteristics of the Venezuelan Oil
  Basins}. In Proceedings of the 3rd World Petroleum Congress, The Hague, the Netherlands,  28 May--6 June 1951 .

\bibitem[Mann \em{et~al.}(2007)Mann, Kroehler, Escalona, Magnani, and
  Christeson]{Mann}
Mann, P.; Kroehler, M.; Escalona, A.; Magnani, B.; Christeson, G. {Subduction
  Along the South Caribbean Deformed Belt: Age of Initiation and Backthrust
  Origin}.
\newblock In {\em {American Geophysical Union}}; {Fall Meeting 2007}; AGU: Washington, DC, USA,
  2007; Chapter Abstract ID. T11D-02, p.~02.

\bibitem[Rodríguez \em{et~al.}(2021)Rodríguez, Bulnes, Poblet, Masini, and
  Flinch]{Rodriguez2021}
Rodríguez, I.; Bulnes, M.; Poblet, J.; Masini, M.; Flinch, J.
\newblock {Structural style and evolution of the offshore portion of the Sinu
  Fold Belt (South Caribbean Deformed Belt) and adjacent part of the Colombian
  Basin}.
\newblock {\em Mar. Pet. Geol.} {\bf 2021}, {\em 125},~104862.
\newblock {{https://doi.org/10.1016/j.marpetgeo.2020.104862}}.

\bibitem[Schubert(1982)]{Schubert}
Schubert, C.
\newblock {Origin of the Cariaco Basin, Southern Caribbean Sea}.
\newblock {\em Mar. Geol.} {\bf 1982}, {\em 47},~345--360.
\newblock {{https://doi.org/10.1016/0025-3227(82)90076-7}}.

\bibitem[Speed \em{et~al.}(1989)Speed, Torrini, and Smith]{Speed}
Speed, R.; Torrini, R.; Smith, P.L.
\newblock {Tectonic evolution of the Tobago Trough forearc basin}.
\newblock {\em J. Geophys. Res.} {\bf 1989}, {\em
  94},~2913--2936.
\newblock {{https://doi.org/10.1029/JB094iB03p02913}}.

\bibitem[Dhont \em{et~al.}(2005)Dhont, Backé, and Hervouët]{Dhont}
Dhont, D.; Backé, G.; Hervouët, Y.
\newblock {Plio-Quaternary extension in the Venezuelan Andes: Mapping from SAR
  JERS imagery}.
\newblock {\em Tectonophysics} {\bf 2005}, {\em 399},~293--312.
\newblock {{https://doi.org/10.1016/j.tecto.2004.12.027}}.

\bibitem[Gibbs and Barron(1993)]{Gibbs}
Gibbs, A.K.; Barron, C.N.
\newblock {\em Geology of the Guiana Shield}; Clarendon Press: Oxford,
  UK, 1993; Volume~22,
\newblock 246p.

\bibitem[Hollowell and Reynolds(2005)]{Hollowell}
Hollowell, T.; Reynolds, R.P. {Checklist of the Terrestrial Vertebrates of the
  Guiana Shield}.
\newblock In {\em {Bulletin of the Biological Society of Washington}};
  Biological Society of Washington: Washington, DC, USA,  2005; Chapter~13, pp. 1--106.

\bibitem[Hammond(2005)]{Hammond}
Hammond, D.S.
\newblock {\em Tropical Forests of the Guiana Shield}; CABI Publishing:
  Wallingford, UK,  2005.

\bibitem[Escalona and Mann(2006)]{Escalona}
Escalona, A.; Mann, P.
\newblock {An overview of the petroleum system of Maracaibo Basin}.
\newblock {\em AAPG Bull.} {\bf 2006}, {\em 90},~657--678.
\newblock {{https://doi.org/10.1306/10140505038}}.

\bibitem[Bockmeulen \em{et~al.}(1983)Bockmeulen, Barker, and
  Dickey]{Bockmeulen}
Bockmeulen, H.; Barker, C.; Dickey, P.A.
\newblock {Geology and geochemistry of crude oils, Bolivar coastal fields,
  Venezuela}.
\newblock {\em AAPG Bull.} {\bf 1983}, {\em 67},~242--270.
\newblock {{https://doi.org/10.1306/03b5acf5-16d1-11d7-8645000102c1865d}}.

\bibitem[Escobar \em{et~al.}(2011)Escobar, Márquez, Inciarte, Rojas, Esteves,
  and Malandrino]{Escobar}
Escobar, M.; Márquez, G.; Inciarte, S.; Rojas, J.; Esteves, I.; Malandrino, G.
\newblock {The organic geochemistry of oil seeps from the Sierra de Perijá
  eastern foothills, Lake Maracaibo Basin, Venezuela}.
\newblock {\em Org. Geochem.} {\bf 2011}, {\em 42},~727--738.
\newblock {{https://doi.org/10.1016/j.orggeochem.2011.06.005}}.

\bibitem[Memon \em{et~al.}(2017)Memon, Borman, Mohammadzadeh, Garcia,
  Tristancho, and Ratulowski]{Memon}
Memon, A.; Borman, C.; Mohammadzadeh, O.; Garcia, M.; Tristancho, D.J.R.;
  Ratulowski, J.
\newblock {Systematic evaluation of asphaltene formation damage of black oil
  reservoir fluid from Lake Maracaibo, Venezuela}.
\newblock {\em Fuel} {\bf 2017}, {\em 206},~258--275.
\newblock {{https://doi.org/10.1016/j.fuel.2017.05.024}}.

\bibitem[Li and Götze(2001)]{Li}
\textls[-15]{Li, X.; Götze, H.J.
\newblock {Ellipsoid, geoid, gravity, geodesy, and geophysics}.
\newblock {\em Geophysics} {\bf 2001}, {\em 66},~1660--1668.
\newblock {{https://doi.org/10.1190/1.1487109}}.}

\bibitem[Yilmaz and Kozlu(2018)]{Yilmaz}
Yilmaz, M.; Kozlu, B.
\newblock {The Comparison of Gravity Anomalies based on Recent High-Degree
  Global Models}.
\newblock {\em Afyon Kocatepe Univ. J. Sci. Eng.}
  {\bf 2018}, {\em 18},~981--990.
\newblock {{https://doi.org/10.5578/fmbd.67502}}.

\bibitem[Guenther(1980)]{Guenther}
Guenther, F.B.
\newblock {Relations between Pasadena magnitude and vertical ground
  acceleration from long distance earthquakes, derived by gravity earthtide
  measurements at Caracas, Venezuela}.
\newblock {\em Phys. Earth Planet. Inter.} {\bf 1980}, {\em
  21},~261--266.
\newblock {{https://doi.org/10.1016/0031-9201(80)90074-6}}.

\bibitem[Schmitz \em{et~al.}(2008)Schmitz, Avila, Bezada, Vieira, Yáñez,
  Levander, Zelt, Jácome, and Magnani]{Schmitz2008}
Schmitz, M.; Avila, J.; Bezada, M.; Vieira, E.; Yáñez, M.; Levander, A.;
  Zelt, C.A.; Jácome, M.I.; Magnani, M.B.
\newblock {Crustal thickness variations in Venezuela from deep seismic
  observations}.
\newblock {\em Tectonophysics} {\bf 2008}, {\em 459},~14--26.
\newblock {{https://doi.org/10.1016/j.tecto.2007.11.072}}.

\bibitem[Ambeh and Lynch(1993)]{Ambeh}
Ambeh, W.B.; Lynch, L.L.
\newblock {Coda Q in the Eastern Caribbean, West Indies}.
\newblock {\em Geophys. J. Int.} {\bf 1993}, {\em
  112},~507--516.
\newblock {\hl{https://doi.org/10.1111/j.1365-246X.1993.tb01184.x}}. %MDPI: Short Biography of Authors is not acceptable in this type of articles, we removed them, please note it.

\end{thebibliography}

%\nocite{*}
 
% \section*{Short Biography of Authors}
%\bio
%{\raisebox{-1.55cm}{\includegraphics[width=3.5cm,height=5.3cm,clip,keepaspectratio]{Photo_author_1.jpg}}}
%{\textbf{Polina Lemenkova} is an Engineer and a PhD student of Prof. Dr. O. Debeir in the Image Analysis Research Unit (LISA-IA) in École polytechnique de Bruxelles, Université Libre de Bruxelles. Her research interests are focused on image processing, cartography, spatial data analysis, remote sensing and applied programming. She has a strong research experience in Earth and environmental sciences, mapping and data visualization, specifically in geological and geophysical domains. Her current topic is on seismic data processing using machine learning methods and scripting languages.}
%%
%\bio
%{\raisebox{-1.9cm}{\includegraphics[width=3.5cm,height=5.3cm,clip,keepaspectratio]{Photo_author_2.jpg}}}
%{\textbf{Olivier Debeir} is a Head of the Image Analysis Research Unit (LISA-IA) in the École polytechnique de Bruxelles (Brussels Faculty of Engineering), Université Libre de Bruxelles. His research interest is in addressing questions related to artificial intelligence and machine learning in general, and~specifically on computer vision, graphics, image analysis, segmentation and classification, pattern recognition and deep learning. His research activity also concerns applied statistics and diverse algorithms of data analysis and modelling in natural sciences. He is an expert in programming languages.}

%%%%%%%%%%%%%%%%%%%%%%%%%%%%%%%%%%%%%%%%%%
\end{adjustwidth}
\end{document}